\documentclass[10pt]{article}

\usepackage{template/response_slu}
\usepackage{booktabs}
\usepackage{array}
\usepackage{amsmath}
\usepackage{longtable}
\usepackage{tabularx}
\usepackage{threeparttable}
\usepackage{float}
\renewcommand{\thetable}{R\arabic{table}}
\RenewDocumentCommand{\checkresponse}{+m}{%
  \par\vspace{0.25\baselineskip}%
  \begingroup
    \color{black}\fontsize{10}{12}\selectfont\justifying
    #1\par
  \endgroup
}

\graphicspath{
  {../../manuscript/figures/round2_candidate/manual/}
  {../../manuscript/figures/round2_candidate/main/}
  {../../manuscript/figures/round2_candidate/appendix/}
  {../../manuscript/figures/}
}

\newcommand{\tablecaptionstyle}[1]{\par\vspace{0.25em}{\footnotesize #1}\par}
\newcommand{\compactfig}[2]{%
  \begin{center}
  \includegraphics[width=#1\linewidth]{#2}
  \end{center}
}
\newcommand{\CompactDescriptorTable}{%
\begin{center}
\footnotesize
\begin{tabularx}{0.96\linewidth}{@{}p{0.15\linewidth}p{0.10\linewidth}p{0.34\linewidth}X@{}}
\toprule
Descriptor & Unit & Calculation & Note \\
\midrule
FAR & -- & $A_\mathrm{gross}/A_\mathrm{block}$ & density \\
SD & m & $h_\mathrm{max}-\bar h$ & height variation \\
AF & floors & $A_\mathrm{gross}/A_\mathrm{foot}$ & linked to FAR, BD \\
$AR_{ew}$ & -- & $\bar h/W_{ew}$ & east-west canyon \\
$AR_{ns}$ & -- & $\bar h/W_{ns}$ & north-south canyon \\
SVF & -- & analytic openness proxy & sky-view openness \\
BD & -- & $A_\mathrm{foot}/A_\mathrm{block}$ & footprint density \\
OSR & -- & $A_\mathrm{open}/A_\mathrm{gross}$ & open-space ratio \\
SC & -- & $A_\mathrm{env}/A_\mathrm{gross}$ & envelope exposure \\
PAR & m$^{-1}$ & $P_\mathrm{sum}/A_\mathrm{foot}$ & compactness \\
$\theta$ & deg & sampled in $[-45,45]$ & orientation \\
OSLI & integer & $0$--$8$ & open-space cell \\
\bottomrule
\end{tabularx}
\end{center}
\tablecaptionstyle{\textbf{Revised Table 1 excerpt.} The 12 quantities are derived morphology descriptors rather than independently sampled design variables.}
}
\newcommand{\EvaluationModeTable}{%
\begin{center}
\footnotesize
\begin{tabular}{llll}
\toprule
Mode & Source & Output & Boundary \\
\midrule
Analytic response generation & generated descriptors & EUIt, EG, H & analytic surface \\
Surrogate prediction & 12 descriptors & DNN outputs & selected checkpoint \\
Feasible projection & descriptor candidate & nearest block & nearest neighbour \\
Physical stress test & locked blocks & physical outputs & 24 cases \\
Climate sensitivity & four blocks & weather deltas & no retraining \\
\bottomrule
\end{tabular}
\end{center}
\tablecaptionstyle{\textbf{Revised Table 2 excerpt.} Evaluation modes are separated by source, output, and interpretation boundary.}
}
\newcommand{\SurrogateMetricsTable}{%
\begin{center}
\footnotesize
\begin{tabular}{lrrrrrr}
\toprule
Validation regime & EUIt nMAE & EG nMAE & H nMAE & EUIt $\rho$ & EG $\rho$ & H $\rho$ \\
\midrule
Repeated 5$\times$5 CV & 0.0176 & 0.0265 & 0.0145 & 0.991 & 0.984 & 0.995 \\
Leave-one-OSLI-out & 0.0185 & 0.0342 & 0.0185 & 0.990 & 0.972 & 0.992 \\
Outer-shell holdout & 0.0278 & 0.0335 & 0.0208 & 0.985 & 0.975 & 0.994 \\
Feature-tail holdout & 0.0230 & 0.0287 & 0.0198 & 0.986 & 0.983 & 0.992 \\
\bottomrule
\end{tabular}
\end{center}
\tablecaptionstyle{\textbf{Revised Table 3 excerpt.} nMAE and Spearman $\rho$ compare DNN predictions with analytic target values.}
}
\newcommand{\BenchmarkTable}{%
\begin{center}
\footnotesize
\begin{tabular}{llrrl}
\toprule
Method & Scenario & HV & IGD & Note \\
\midrule
DDPG & Balanced & 0.661912 & 0.465522 & scalarized output \\
DDPG & Saving & 1.002757 & 0.276529 & scalarized output \\
DDPG & Generation & 1.170428 & 0.097551 & scalarized output \\
NSGA-II & equal-size 20 & 1.330995 & 0.004710 & common retained size \\
NSGA-II & full archive & 1.330999 & 0.004947 & 2000 rows \\
\bottomrule
\end{tabular}
\end{center}
\tablecaptionstyle{\textbf{Revised Table 5 excerpt.} Descriptor-space metrics use a fixed normalized reference front.}
}
\newcommand{\PhysicalSummaryTable}{%
\begin{center}
\footnotesize
\begin{tabular}{lrrrl}
\toprule
Target & Count & MAE & nMAE & Spearman $\rho$ \\
\midrule
EUIt & 18 & 3.788 & 0.340 & -0.404 \\
EG GHI proxy & 18 & 0.441 & 0.405 & -0.023 \\
H & 18 & 1.108 & 0.601 & -0.043 \\
\bottomrule
\end{tabular}
\end{center}
\tablecaptionstyle{\textbf{Revised Table 6 and Table S4 excerpt.} The direct feasible stress-test cases show limited metric and rank consistency.}
}
\newcommand{\ClimateTable}{%
\begin{center}
\footnotesize
\begin{tabular}{lrrrl}
\toprule
Station & $\Delta$EUIt & $\Delta$EG & $\Delta$H & Main interpretation \\
\midrule
Beijing & 13.467 & 0.135 & -0.056 & higher demand \\
Guangzhou & -2.972 & -0.057 & -0.025 & lower demand \\
Harbin & 55.681 & -0.044 & -0.675 & strong demand increase \\
\bottomrule
\end{tabular}
\end{center}
\tablecaptionstyle{\textbf{Revised Table S5 excerpt.} $\Delta$EUIt is in kWh m$^{-2}$ y$^{-1}$, $\Delta$EG is in $10^6$ kWh y$^{-1}$, and $\Delta$H is in h.}
}
\newcommand{\RewardEquations}{%
\[
z_i =
\frac{y_i-y_i^{\min}}
{y_i^{\max}-y_i^{\min}} .
\]
\[
\mathbf{u} = [0,1,1]^\mathsf{T}.
\]
\[
d_w =
\frac{\left\|\mathbf{w}\odot(\mathbf{z}-\mathbf{u})\right\|_2}
{\left\|\mathbf{w}\right\|_2},
\qquad
R = 1-d_w .
\]
}
\newcommand{\AnalyticEquations}{%
\[
p_\theta = |\theta| / 45,\qquad
a_m = 0.5(AR_{ew}+AR_{ns}),
\]
\[
s = \exp\left(-\frac{(FAR-1.7)^2}{0.6}\right)(0.45+SVF)(0.3+OSR),
\qquad
p_d = \max(FAR-2.8,0).
\]
\[
\begin{aligned}
EUIt^\ast ={}& 78.0 + 10.0BD + 4.3PAR + 2.8a_m + 0.022SD + 0.8SC + 2.1p_\theta \\
& - 6.4OSR - 5.8SVF - 0.9AF - 4.5s + \epsilon_{EUIt},\\
EG^\ast ={}& 0.96 + 0.28FAR + 0.64SVF + 0.52OSR - 0.19BD - 0.12a_m - 0.18p_\theta \\
& + 0.26OSLI^\ast + 0.78s - 0.1p_d + \epsilon_{EG},\\
H^\ast ={}& 6.45 + 1.1SVF + 0.55OSR + 0.48OSLI^\ast + 0.24s - 0.7p_\theta - 0.45PAR - 0.24a_m + \epsilon_H .
\end{aligned}
\]
\[
EUIt = \operatorname{clip}(EUIt^\ast,66.0,96.0),\quad
EG = \operatorname{clip}(EG^\ast,1.2,2.85),\quad
H = \operatorname{clip}(H^\ast,6.0,7.85).
\]
}

\renewcommand{\responsetitletext}{Response to Reviewers}
\renewcommand{\responseopening}{%
Dear Editor,\par
\vspace{0.7em}
We sincerely thank the editor and reviewers for their careful reading of our manuscript and for the constructive comments. We have revised the manuscript, Supplementary Information, figures, tables, equations, references, and code/data availability statement to make the data-generation workflow, optimizer formulation, comparison protocol, physical stress-test boundary, and climate-sensitivity boundary explicit.\par
\vspace{0.7em}
This response letter has been reorganized as a comment-level response package. Each reviewer comment is reproduced under its own comment heading. The response then states the revision, identifies where the manuscript was changed, and displays the relevant revised manuscript, Supplementary Information, figure, table, equation, or local reference material directly under that comment.\par
\vspace{0.7em}
\noindent\textbf{Manuscript title:} Surrogate-conditioned optimizer ranking in block-scale urban energy design: A comparison of policy learning and Pareto search\par
\noindent\textbf{Manuscript ID:} APEN-D-26-08148\par
\vspace{0.7em}
Sincerely,\par
The Authors%
}

\begin{document}
\makeresponsetitle
\clearpage
\responsecontents

\section*{Summary of major revisions}
\addcontentsline{toc}{section}{Summary of major revisions}

\responsepara{%
We are grateful to the editor and reviewers for the careful and constructive evaluation of our manuscript. After rereading the original submission in light of the comments, we recognize that several parts of the previous version were not sufficiently clear, especially the data-generation workflow, the role of the analytic response surface, the DDPG formulation, the optimizer comparison protocol, and the interpretation of the physics-based and climate-sensitive calculations. We apologize for the additional effort that these ambiguities caused during review. The comments were highly valuable and led us to revise the manuscript, figures, tables, equations, Supplementary Information, and reproducibility statement in a much more transparent form.
}

\responsepara{%
The revised manuscript is therefore not a minor language edit. We substantially rewrote the study framing, clarified the distinction between generated feasible morphologies and morphology descriptors, replaced ambiguous validation language with a limited physics-based cross-model stress test, added cross-climate sensitivity calculations, corrected the reward formulation, and reorganized the Supplementary Information so that the equations, validation protocols, optimizer analyses, physical case results, and climate case results are directly traceable.
}

\begin{table}[H]
\centering
\caption{Editor-facing overview of the main revision packages.}
{\scriptsize
\setlength{\tabcolsep}{2.4pt}
\renewcommand{\arraystretch}{0.94}
\begin{tabularx}{\linewidth}{@{}p{0.18\linewidth}p{0.18\linewidth}p{0.30\linewidth}X@{}}
\toprule
Revision package & Comments addressed & Manuscript revision & Supporting material \\
\midrule
Framing and readability & R1-1, R1-2, R1-15, R4-1 & Abstract, Introduction, Nomenclature, Discussion, Conclusions & SI S1--S6 aligned with revised terminology \\
Data source and objective definitions & R2-1, R2-6, R2-7, R3-5 & Analytic response generation, evaluation modes, target definitions & SI S1 equations, clipping ranges, descriptor dependencies \\
Sampling and surrogate assessment & R1-7, R1-8, R2-4 & Sampling protocol, descriptor coverage, repeated validation & SI S1--S2, Table S2, Fig.~S1--S3 \\
DDPG and optimizer comparison & R1-10, R1-11, R1-13, R2-5, R3-2, R3-3, R4-4 & Static surrogate-query DDPG, reward correction, output contracts & SI S3, equal-size metrics, HV ceiling analysis \\
Feasible projection and physical stress test & R1-7, R2-2, R2-3, R3-5 & Nearest-feasible-block projection and 24-case cross-model stress test & SI S5, projection metrics, physical summary table \\
Cross-climate and reproducibility & R1-14, R2-8, R4-2, R4-3, R4-5 & Four-block climate sensitivity, future energy-system scope, data/code statement & SI S6, weather metadata, public repository statement \\
\bottomrule
\end{tabularx}
}
\end{table}

\responsepara{%
Table R1 is only an editor-facing overview. The detailed revised text, figures, tables, equations, Supplementary Information excerpts, and local references are provided under the corresponding reviewer comments below.
}

\reviewerpage{1}
\checkresponse{%
We thank Reviewer 1 for the detailed comments on readability, terminology, figures, software context, parameter selection, validation, cross-validation, framework structure, DDPG formulation, units, case generality, and broader-audience accessibility. The revised manuscript responds by simplifying the framing, adding definitions and tables, expanding the method explanation, and keeping the claims within the analytic response-generation and limited sensitivity-analysis scope.
}

\comment{R1-1}
\begin{commentbox}
In the Abstract, the paragraph is too long (more than 250 words), which can be simplified. Furthermore, it is recommended that authors point out the significance of the work with more brief but powerful expressions.
\end{commentbox}
\responseheading{Response.}
\checkresponse{%
We thank the reviewer for this helpful suggestion. We agree that the original Abstract was too long and did not express the significance of the study clearly enough. In the revised manuscript, we rewrote the Abstract so that readers can see the main research question immediately: the work compares optimizer behaviour under a shared surrogate response surface and then examines how this comparison changes after feasible-block projection, selected physics-based calculations, and cross-climate sensitivity analysis.

The revised Abstract now follows a clearer sequence. It first introduces the block-scale urban energy design problem. It then explains that feasible block layouts are generated and converted into 12 morphology descriptors, and that EUIt, EG, and H are assigned by a morphology-to-performance response surface. After that, it reports the selected 2000-sample DNN surrogate, repeated-validation accuracy, the DDPG and NSGA-II comparison, the compression of the NSGA-II archive from 2000 descriptor candidates to 51 unique feasible blocks, the 24-case cross-model stress test, and the four-block cross-climate sensitivity analysis.

We also revised the significance statement to avoid an overly broad optimizer-superiority interpretation. The Abstract now presents the study as a representation-aware design-support workflow. This wording makes the contribution more direct for Applied Energy readers while still preserving the interpretation limits shown by the physical and climate-sensitive calculations.
}
\responseheading{Revisions made in the manuscript.}
\responsepara{%
We made three corresponding revisions. First, we rewrote the Abstract around the problem, method, key numerical results, and interpretation limits. Second, we revised the final Introduction contribution paragraph so that the significance is framed as representation-aware design support. Third, we revised the Conclusions so that the final message emphasizes projection distance, model agreement, and climate sensitivity rather than direct physical optimizer superiority.
}
\responseheading{Relevant revised manuscript and supporting material.}
\begin{revisionbox}
\revisionloc{Main manuscript, Abstract.}
Block-scale urban energy design requires optimization methods that connect urban form with building energy demand, rooftop photovoltaic potential, and winter solar access during early-stage planning. This study develops a surrogate-conditioned framework to compare policy learning and Pareto search for residential-block morphology optimization. Feasible block layouts are generated first and converted into 12 morphology descriptors. A morphology-to-performance response surface then assigns three targets to each sample: Energy Use Intensity, Energy Generation, and winter sunlight hours. These targets are denoted EUIt, EG, and H throughout the paper. The 2000-sample DNN surrogate selected from nested 500/1000/1500/2000-sample pools provides the common environment for DDPG and NSGA-II. Repeated five-fold validation gives mean nMAE values of 0.0176, 0.0265, and 0.0145 for EUIt, EG, and H. Under the shared surrogate and matched query budget, NSGA-II attains near-ceiling descriptor-space HV, while DDPG shows preference-dependent performance. Projection to generated feasible blocks compresses the 2000-candidate NSGA-II archive to 51 unique feasible blocks, indicating that descriptor-space coverage and feasible-morphology diversity differ substantially. A 24-case cross-model stress test completes all locked cases and shows limited metric and rank consistency between analytic and physics-based outputs. A four-block cross-climate sensitivity analysis shows stable EG ordering but climate-dependent EUIt and H responses. The study contributes a representation-aware urban design workflow that links surrogate search, feasible-block projection, and sensitivity analysis to improve the interpretability of optimizer comparisons for early-stage urban energy planning.
\end{revisionbox}

\comment{R1-2}
\begin{commentbox}
Overall, there are too many acronyms which decreases the paper's readability. Moreover, symbols should be added.
\end{commentbox}
\responseheading{Response.}
\checkresponse{%
We thank the reviewer for pointing this out. We agree that the previous front-matter list addressed abbreviations but did not sufficiently support the mathematical notation used in the methodology and optimizer equations. In response, we changed the front-matter block from Abbreviations to Nomenclature and divided it into two parts.

The first part defines the abbreviations used throughout the paper, including EUIt, EG, H, DNN, DDPG, NSGA-II, CMA-ES, GHI, HV, IGD, nMAE, and OSLI. The second part defines the symbols used in the analytic response-generation and DDPG reward formulations, including the morphology descriptor vector, analytic target vector, normalized target coordinate, target bounds, ideal vector, axis-scaling vector, weighted distance, training reward, orientation angle, noise term, and geometric quantities used by the descriptor definitions.

We also simplified Table 1. The previous version repeated the same content in the Descriptor and Symbol columns. The revised Table 1 now keeps only the descriptor name, unit, calculation, and note. This separation is clearer: the Nomenclature supports equation readability, while Table 1 explains how the morphology descriptors are constructed and interpreted.
}
\responseheading{Revisions made in the manuscript.}
\responsepara{%
We made three corresponding revisions. First, we replaced the front-matter Abbreviations block with a Nomenclature block. Second, we added a Formula symbols subsection for the analytic-response and DDPG reward notation. Third, we simplified Table 1 by removing the redundant Symbol column and retaining descriptor, unit, calculation, and note as the compact descriptor summary.
}
\responseheading{Relevant revised manuscript and supporting material.}
\begin{revisionbox}
\revisionloc{Main manuscript, Nomenclature excerpt.}
\begin{center}
\footnotesize
\begin{tabularx}{0.96\linewidth}{@{}p{0.14\linewidth}X p{0.13\linewidth}X@{}}
\toprule
\multicolumn{4}{@{}l}{\textit{Abbreviations}}\\
EUIt & Energy Use Intensity & EG & Energy Generation \\
H & winter sunlight hours & DNN & deep neural network \\
DDPG & Deep Deterministic Policy Gradient & NSGA-II & Non-dominated Sorting Genetic Algorithm II \\
CMA-ES & Covariance Matrix Adaptation Evolution Strategy & GHI & global horizontal irradiance \\
HV & Hypervolume & IGD & Inverted Generational Distance \\
nMAE & normalized mean absolute error & OSLI & open-space location index \\
\bottomrule
\end{tabularx}
\end{center}

\begin{center}
\footnotesize
\begin{tabularx}{0.96\linewidth}{@{}p{0.11\linewidth}X p{0.11\linewidth}X@{}}
\toprule
\multicolumn{4}{@{}l}{\textit{Formula symbols}}\\
$\mathbf{x}$ & morphology descriptor vector & $\mathbf{y}^{ana}$ & analytic target vector \\
$y_i$ & target value for objective $i$ & $z_i$ & normalized target coordinate \\
$y_i^{\min}$ & lower target bound & $y_i^{\max}$ & upper target bound \\
$\mathbf{u}$ & ideal vector $[0,1,1]^\mathsf{T}$ & $\mathbf{w}$ & axis-scaling coefficient vector \\
$d_w$ & weighted normalized distance & $R$ & DDPG training reward \\
$\theta$ & block orientation angle & $\epsilon$ & response-generation noise term \\
$A_{gross}$ & gross floor area & $A_{block}$ & block area \\
$A_{foot}$ & building footprint area & $A_{open}$ & open-space area \\
$A_{env}$ & envelope area & $P_{sum}$ & summed footprint perimeter \\
$\bar{h}$ & mean building height & $h_{\max}$ & maximum building height \\
$W_{ew}$ & east-west canyon width & $W_{ns}$ & north-south canyon width \\
\bottomrule
\end{tabularx}
\end{center}
\revisionseparator
\revisionloc{Main manuscript, Table 1 excerpt.}
\CompactDescriptorTable
\end{revisionbox}

\comment{R1-3}
\begin{commentbox}
Figure 1: The caption should be improved to explain the content of the figure in more detail.
\end{commentbox}
\responseheading{Response.}
\checkresponse{%
We thank the reviewer for this helpful suggestion. We agree that the original Fig. 1 caption was too brief for a workflow figure that introduces the full methodology. In the revised manuscript, Fig. 1 is used to orient the reader before the detailed method sections, so the caption now explains the main computational stages rather than only naming the figure.

We revised both the paragraph before Fig. 1 and the figure caption. The paragraph now explains that the study distinguishes generated feasible block geometries, morphology descriptors used by the surrogate and optimizers, and locked block cases used in selected physics-based and climate-sensitive calculations. The caption then describes the workflow in order: feasible morphology generation, analytic response and surrogate modelling, surrogate-assisted optimization, and feasible design assessment.

This revision also prevents a possible misunderstanding in the original manuscript. The figure and caption now make clear that the 2000-sample benchmark uses analytic response assignment, whereas EnergyPlus- and Radiance-component calculations are used later for selected locked cases in the cross-model stress test and climate sensitivity analysis.

We made this distinction in the caption itself because readers often use Fig. 1 as the entry point to the method. The revised caption therefore carries the workflow boundary before readers reach the more detailed equations and result sections.
}
\responseheading{Revisions made in the manuscript.}
\responsepara{%
We made two corresponding revisions. First, we rewrote the methodology transition before Fig. 1 so that the reader sees the representation distinction before the diagram. Second, we expanded the Fig. 1 caption to explain the four workflow modules and to distinguish analytic response generation from selected physical calculations.
}
\responseheading{Relevant revised manuscript and supporting material.}
\begin{revisionbox}
\revisionloc{Main manuscript, methodology transition before Fig.~1.}
The methodology links morphology generation, analytic target assignment, surrogate-assisted optimization, and post-optimization assessment within one block-scale design workflow. The workflow distinguishes three representations: generated feasible block geometries, 12 morphology descriptors used by the surrogate and optimizers, and locked block cases examined through physics-based and climate-sensitive calculations. This distinction is essential because the optimizer operates in descriptor space, whereas design interpretation ultimately depends on generated block geometries. Figure 1 summarizes the computational sequence.
\revisionseparator
\revisionloc{Main manuscript, Fig.~1.}
\compactfig{0.90}{fig1.pdf}
\figcaptionstyle{\textbf{Revised Fig.~1.} Workflow of the surrogate-assisted block-scale morphology optimization framework. Feasible block layouts are generated and converted into 12 morphology descriptors. The analytic response and surrogate-modelling stage assigns EUIt, EG, and H and trains the selected DNN surrogate. DDPG and NSGA-II are compared under the shared surrogate, then descriptor-space candidates are projected to feasible blocks for selected physical and climate-sensitivity checks.}
\end{revisionbox}

\comment{R1-4}
\begin{commentbox}
Line 36-38: According to the Chineses ~ requires reference. Moreover, the authors need to specify some references for the Grasshopper platform. Please describe the general context of the softwares (tools) and why that softwares are suitable for your analysis.
\end{commentbox}
\responseheading{Response.}
\checkresponse{%
We thank the reviewer for this helpful comment. We agree that the manuscript needed a clearer explanation of both the Chinese urban energy-planning context and the software workflow used for the selected physical calculations. The revised manuscript now separates these two issues so that readers can see why the study is relevant and what role each computational tool plays.

For the planning context, we expanded the Introduction to discuss block-scale energy demand, rooftop PV potential, winter solar access, Chinese urban renewal, and climate-resilient planning. We also added peer-reviewed urban-energy and surrogate-optimization studies to place the work within existing Applied Energy and urban-building literature.

For the software context, we revised the method text to explain that the selected physics-based and climate-sensitive cases are implemented through a Grasshopper/Rhino-based workflow linking generated block geometry to Honeybee, EnergyPlus, and Radiance components. This workflow is used only for locked case calculations in the cross-model stress test and cross-climate sensitivity analysis. The 2000-sample optimizer benchmark remains an analytic response-generation dataset, which provides a controlled response surface for comparing DDPG and NSGA-II under repeated surrogate queries.

This wording is intentionally bounded. We do not claim that all 2000 samples are direct Grasshopper/EnergyPlus/Radiance simulations, and we do not present the selected physical workflow as full validation of the analytic benchmark. Instead, the revised manuscript explains how the selected physical calculations are connected to the broader surrogate-assisted benchmark.
}
\responseheading{Revisions made in the manuscript.}
\responsepara{%
We made three corresponding revisions. First, we expanded the Introduction around Chinese urban renewal, PV-oriented planning, and climate resilience. Second, we revised Section 2.6 to explain the Grasshopper/Rhino, Honeybee, EnergyPlus, and Radiance workflow for selected locked cases. Third, we separated this selected-case workflow from the analytic response-generation benchmark used by the optimizers.
}
\responseheading{Relevant revised manuscript and supporting material.}
\begin{revisionbox}
\revisionloc{Main manuscript, Motivation paragraph.}
Block morphology influences energy performance through coupled physical mechanisms. Density and height alter heating and cooling demand by changing envelope exposure, mutual shading, and microclimatic context. Solar access and daylight availability depend on roof exposure, sky-view conditions, block orientation, and open-space arrangement. Recent large-scale studies further show that urban spatial structure can shape building energy demand across broad urban areas and under extreme climate events, which makes morphology-sensitive design relevant beyond isolated building cases.
\revisionseparator
\revisionloc{Main manuscript, Section 2.2 paragraph.}
This module provides a controlled morphology-to-performance response surface for optimizer comparison. The nested 500/1000/1500/2000-sample pools are deterministic subsets of the same generated morphology pool under the recorded random seed and morphology-generation protocol. The physics-based case calculations are conducted subsequently on locked generated blocks to assess how the analytic target values relate to EnergyPlus- and Radiance-component outputs in selected cases.
\revisionseparator
\revisionloc{Main manuscript, Section 2.6 software-workflow paragraph.}
The selected physics-based and climate-sensitive cases are implemented through a Grasshopper/Rhino-based workflow that links generated block geometry to Honeybee, EnergyPlus, and Radiance components. This workflow is used only for the locked case calculations in the physics-based cross-model stress test and cross-climate sensitivity analysis. The 2000-sample optimizer benchmark remains an analytic response-generation dataset, which provides the controlled response surface for comparing DDPG and NSGA-II under repeated surrogate queries.
\revisionseparator
\revisionloc{Supplementary Material, Section S5.3 paragraph.}
The 18 direct feasible cases completed both EnergyPlus and Radiance-component runs. EUIt agreement is weak, EG is represented by a GHI-based rooftop-PV proxy, and H is represented by January 20 windowsill direct-sun hours. The nMAE and rank-correlation values therefore define a stress-test boundary for the tested direct cases.
\end{revisionbox}
\responseheading{Local references.}
\begin{responsereferences}
\referenceitem{Javanroodi, K., Perera, A. T. D., et al. (2023). Designing climate resilient energy systems in complex urban areas considering urban morphology: A technical review. \textit{Advances in Applied Energy}, 12. \responsedoi{https://doi.org/10.1016/j.adapen.2023.100155}}
\referenceitem{Du, L., Wang, H., et al. (2025). Impact of block form on building energy consumption, urban microclimate and solar potential: A case study of Wuhan, China. \textit{Energy and Buildings}, 328. \responsedoi{https://doi.org/10.1016/j.enbuild.2024.115224}}
\referenceitem{Wang, H., Wu, Z., Gu, W., et al. (2026). Towards climate-resilient cities: Assessing and planning urban energy demand under extreme climate events. \textit{Energy and Buildings}, 337, 117436. \responsedoi{https://doi.org/10.1016/j.enbuild.2026.117436}}
\end{responsereferences}

\comment{R1-5}
\begin{commentbox}
Line 49: It would be good to describe why you select 12 parameters for building models.
\end{commentbox}
\responseheading{Response.}
\checkresponse{%
We thank the reviewer for this comment. We agree that the descriptor-selection logic and dependency structure were not explained clearly enough. This point matters because the 12 morphology quantities are derived descriptors from generated block geometries rather than independent variables that can vary freely.

We revised Section 2.1, Table 1, and Supplementary Sections S1.2--S1.5 to rewrite the descriptor explanation, define the descriptors with units and construction notes, and add supporting material on prototype settings, OSLI, dependency, PCA coverage, and distributions. The revised material now gives the reader a more direct account of what was changed and how the revision should be interpreted.

The added calculations show that FAR, BD, AF, and OSR share footprint and gross-floor-area quantities, that OSLI is integer-valued by construction, that six principal components explain 95 percent of descriptor variance, and that exact duplicate descriptor rows are absent in the canonical 2000-row pool. Optimizer outputs are therefore interpreted first in descriptor space and then through nearest-feasible-block projection before morphology-level conclusions are drawn.
}
\responseheading{Revisions made in the manuscript.}
\responsepara{%
We made three corresponding revisions. First, we rewrote the descriptor-selection paragraph around density, height, openness, envelope exposure, compactness, orientation, and open-space position. Second, we revised Table 1 to include unit or type, construction, and interpretation. Third, we added supporting material on descriptor dependency and PCA coverage.
}
\responseheading{Relevant revised manuscript and supporting material.}
\begin{revisionbox}
\revisionloc{Main manuscript, Section 2.1 paragraphs.}
The morphology generator first constructs feasible block layouts and only then calculates descriptor values. Each block is represented as a 240 m $\times$ 240 m square divided into a $3\times3$ grid of 80 m $\times$ 80 m land units. One cell is assigned as open space, and the remaining eight cells are populated with residential prototypes. For each generated block, the generator samples the open-space position, building prototype, number of floors, and orientation angle. The orientation angle is sampled within -45 deg to 45 deg. The resulting block is therefore a generated feasible morphology, while the 12 input quantities used by the surrogate are descriptors derived from that morphology.

The descriptor set in Table 1 includes FAR, SD, AF, AR\_ew, AR\_ns, SVF, BD, OSR, SC, PAR, theta, and OSLI. These quantities are derived morphology descriptors with geometric dependencies. For example, FAR, BD, AF, and OSR are linked through the same footprint and gross-floor-area quantities, while OSLI is an integer descriptor of open-space position. Consequently, optimizer outputs are first interpreted in descriptor space and then mapped to generated feasible blocks before morphology-level conclusions are drawn.
\revisionseparator
\revisionloc{Main manuscript, Table 1 excerpt.}
\CompactDescriptorTable
\revisionseparator
\revisionloc{Supplementary Material, Table S2 excerpt.}
\begin{center}
\footnotesize
\begin{tabular}{lll}
\toprule
Check & Value & Interpretation \\
\midrule
FAR minus BD times AF & $1.33\times10^{-15}$ & numerical roundoff \\
OSR relation residual & $2.22\times10^{-16}$ & footprint relation holds \\
OSLI integer consistency & 1.0 & all OSLI values are integer \\
Exact duplicate count & 0 & no exact duplicate descriptor rows \\
PCA components for 95\% variance & 6 & descriptor space is structured \\
\bottomrule
\end{tabular}
\end{center}
\revisionseparator
\revisionloc{Supplementary Material, Fig.~S1.}
\compactfig{0.86}{A1_descriptor_distributions.pdf}
\figcaptionstyle{\textbf{Fig.~S1 excerpt.} Descriptor distributions and nearest-neighbour distances in the generated morphology pool.}
\end{revisionbox}

\comment{R1-6}
\begin{commentbox}
Table 1: The author could add an additional column on the right to describe Factors (Floor area ratio, Setback, diversity, etc.). Table 2: The parameters are not assumptions; it seems the baseline or operating parameters for the building energy models.
\end{commentbox}
\responseheading{Response.}
\checkresponse{%
We thank the reviewer for this comment. We agree that the functions of Table 1 and Table 2 were not sufficiently distinct. This point matters because readers need to know which table defines morphology inputs and which table defines evaluation modes.

We revised Tables 1 and 2 and Sections 2.1--2.2 to separate descriptor construction from evaluation-mode interpretation. The revised material now gives the reader a more direct account of what was changed and how the revision should be interpreted.

Table 1 now defines morphology descriptors by unit or type, construction, and interpretation, while Table 2 distinguishes analytic response generation, surrogate prediction, feasible projection, physical stress testing, and climate sensitivity by source, output, and boundary. This split avoids mixing descriptor construction, target assignment, surrogate prediction, and selected physical calculations in one undifferentiated method description.

This table separation also makes later revisions easier to follow because response-surface construction, surrogate prediction, and physical stress testing now have different table anchors.
}
\responseheading{Revisions made in the manuscript.}
\responsepara{%
We made three corresponding revisions. First, we rewrote Table 1 as the descriptor-definition table. Second, we rewrote Table 2 as the evaluation-mode table. Third, we aligned the surrounding method text with the two-table division.
}
\responseheading{Relevant revised manuscript and supporting material.}
\begin{revisionbox}
\revisionloc{Main manuscript, Table 1 excerpt.}
\CompactDescriptorTable
\revisionseparator
\revisionloc{Main manuscript, Table 2 excerpt.}
\EvaluationModeTable
\end{revisionbox}

\comment{R1-7}
\begin{commentbox}
The author may need to justify the methodology used. I think it would be great to have model validation studies using published literature data or exp. data. Moreover, the authors need to describe the underlying assumptions of their models.
\end{commentbox}
\responseheading{Response.}
\checkresponse{%
We thank the reviewer for this important suggestion. We agree that the original manuscript did not make the modelling assumptions and validation hierarchy sufficiently clear. We also agree that readers should not be asked to accept surrogate accuracy as if it were direct physical or experimental validation.

In the revised manuscript, we therefore separate three levels of assessment. First, the DNN surrogate is validated against the analytic response-generation targets through repeated five-fold cross-validation and boundary-oriented holdout checks. Second, the analytic response surface is documented explicitly through formulas, stochastic terms, clipping ranges, and descriptor dependencies. Third, selected generated blocks are examined through a limited physics-based cross-model stress test using EnergyPlus- and Radiance-component calculations.

The revised results are deliberately conservative. The repeated 5 by 5 surrogate validation gives nMAE values of 0.0176, 0.0265, and 0.0145 for EUIt, EG, and H, respectively, but this only measures agreement with analytic targets. The physical stress test gives weaker agreement for the 18 direct feasible cases, with Spearman \(\rho\) values of -0.404 for EUIt, -0.023 for the EG proxy, and -0.043 for H. We added related peer-reviewed studies to motivate the morphology-energy relationship, but we do not treat those studies as direct validation data because the datasets and model protocols differ.

This revision directly narrows the claim. The physical calculations now serve as a stress test showing where analytic and physics-based outputs diverge, not as proof that the optimizer ranking transfers to a fully physical model.
}
\responseheading{Revisions made in the manuscript.}
\responsepara{%
We made four corresponding revisions. First, we added the analytic-response equations, stochastic terms, and clipping ranges to the supporting material. Second, we expanded the surrogate-validation section to report repeated and boundary-oriented checks. Third, we added the 24-case physical stress-test summary. Fourth, we rewrote the Discussion so that weak physical agreement limits descriptor-space optimizer interpretation.
}
\responseheading{Relevant revised manuscript and supporting material.}
\begin{revisionbox}
\revisionloc{Main manuscript, Section 2.2 paragraph.}
In this study, the analytic response generator is the target-assignment module that maps a generated morphology descriptor vector, denoted $\mathbf{x}$, to three target responses, denoted $\mathbf{y}^{ana}=[\mathrm{EUIt},\mathrm{EG},H]^\mathsf{T}$. The mapping uses explicit formulas that combine density, openness, aspect ratio, sky-view factor, orientation, and open-space position to represent demand, rooftop-generation, and sunlight-access responses. A small stochastic term is added during dataset construction, and the target values are then limited to the response ranges used for surrogate training.
\revisionseparator
\revisionloc{Main manuscript, Table 3 excerpt.}
\SurrogateMetricsTable
\revisionseparator
\revisionloc{Main manuscript, Fig.~9.}
\compactfig{0.86}{physical_cross_model_stress_test.pdf}
\figcaptionstyle{\textbf{Revised Fig.~9.} Cross-model stress test for 18 direct generated feasible blocks.}
\revisionseparator
\revisionloc{Supplementary Material, Table S4 excerpt.}
\PhysicalSummaryTable
\end{revisionbox}

\comment{R1-8}
\begin{commentbox}
The author may need to describe what the five-fold cross validation means.
\end{commentbox}
\responseheading{Response.}
\checkresponse{%
We thank the reviewer for this comment. We agree that the cross-validation explanation should be understandable beyond the metric values. This point matters because small error numbers are not enough unless the reader knows what was held out and what source layer was being tested.

We revised Section 2.3, Table 3, Fig. 5, and Supplementary Sections S2.2--S2.5 to explain repeated five-fold cross-validation as repeated held-out testing of the selected 2000-sample DNN surrogate against analytic target values. The revised material now gives the reader a more direct account of what was changed and how the revision should be interpreted.

The revised validation text keeps the repeated nMAE and Spearman metrics and adds OSLI holdout, outer-shell holdout, and feature-tail holdout checks for more difficult descriptor regions. The assessment compares DNN agreement with analytic targets, not agreement with EnergyPlus or Radiance outputs.
}
\responseheading{Revisions made in the manuscript.}
\responsepara{%
We made three corresponding revisions. First, we rewrote the cross-validation explanation in plain methodological terms. Second, we retained the repeated 5 by 5 nMAE and Spearman values in Table 3. Third, we added boundary-oriented validation families in the supporting material.
}
\responseheading{Relevant revised manuscript and supporting material.}
\begin{revisionbox}
\revisionloc{Main manuscript, Section 2.3 paragraph.}
The selected surrogate is assessed through repeated five-fold cross-validation, leave-one-OSLI-out validation, outer-shell holdout, and feature-tail holdout. The repeated five-fold protocol evaluates interpolation performance under randomized folds. The leave-one-OSLI-out protocol tests whether the surrogate remains stable when one open-space-position class is withheld. The outer-shell and feature-tail holdouts examine surrogate behaviour near the boundary of the sampled descriptor domain.
\revisionseparator
\revisionloc{Main manuscript, Table 3 excerpt.}
\SurrogateMetricsTable
\revisionseparator
\revisionloc{Main manuscript, Fig.~5.}
\compactfig{0.86}{surrogate_robustness.pdf}
\figcaptionstyle{\textbf{Revised Fig.~5.} Surrogate assessment across random, grouped, and boundary-oriented validation regimes.}
\revisionseparator
\revisionloc{Supplementary Material, Section S2.2 paragraph.}
Repeated five-fold cross-validation gives mean nMAE values of 0.0176 for EUIt, 0.0265 for EG, and 0.0145 for H. The corresponding Spearman correlations are 0.991, 0.984, and 0.995. These values support low error within the analytic response-generation regime.
\end{revisionbox}

\comment{R1-9}
\begin{commentbox}
The framework has two components: The first .... The second ...
\end{commentbox}
\responseheading{Response.}
\checkresponse{%
We thank the reviewer for this comment. We agree that the framework components needed to be introduced before the detailed results. This point matters because the workflow is easier to audit when each result is tied to the representation and evaluation mode that produced it.

We revised the Methodology opening, Fig. 1, Table 2, and Section 2.6 to rewrite the workflow as a sequence of feasible morphology generation, descriptor extraction, analytic response assignment, surrogate training, optimizer querying, projection, and selected sensitivity analyses. The revised material now gives the reader a more direct account of what was changed and how the revision should be interpreted.

The revised workflow lets the reader trace whether a result comes from the analytic benchmark, the selected DNN surrogate, feasible projection, physical stress testing, or climate sensitivity. The full 2000-sample benchmark is now explicitly described as analytic response assignment plus surrogate prediction, not a set of direct annual physical simulations.
}
\responseheading{Revisions made in the manuscript.}
\responsepara{%
We made three corresponding revisions. First, we rewrote the opening methodology paragraph around workflow components. Second, we revised Fig. 1 and its caption to show the same component order. Third, we used Table 2 to define each evaluation mode and its boundary.
}
\responseheading{Relevant revised manuscript and supporting material.}
\begin{revisionbox}
\revisionloc{Main manuscript, methodology opening paragraph.}
The methodology links morphology generation, analytic target assignment, surrogate-assisted optimization, and post-optimization assessment within one block-scale design workflow. The workflow distinguishes three representations: generated feasible block geometries, 12 morphology descriptors used by the surrogate and optimizers, and locked block cases examined through physics-based and climate-sensitive calculations.
\revisionseparator
\revisionloc{Main manuscript, Fig.~1.}
\compactfig{0.90}{fig1.pdf}
\figcaptionstyle{\textbf{Revised Fig.~1.} Workflow linking feasible morphology generation, analytic response assignment, surrogate-assisted optimization, feasible-block projection, and sensitivity analysis.}
\revisionseparator
\revisionloc{Main manuscript, Table 2 excerpt.}
\EvaluationModeTable
\end{revisionbox}

\comment{R1-10}
\begin{commentbox}
Please add a reference for Markov Decision Process.
\end{commentbox}
\responseheading{Response.}
\checkresponse{%
We thank the reviewer for pointing this out. We added foundational references for Markov Decision Processes and reinforcement learning, and we also revised the surrounding text to avoid overstating the MDP interpretation of the morphology problem.

The revised manuscript now explains that MDPs provide the general reinforcement-learning formalism for sequential decision-making, while the present DDPG implementation should be read more narrowly. In this study, the environment serializes repeated queries to a static surrogate response surface. The actor outputs an absolute 12-descriptor vector, the surrogate returns the next normalized target state, and one episode is a fixed sequence of 40 surrogate queries. This is different from a physical-time urban morphology control process.

This revision addresses the requested citation and clarifies the modelling boundary at the same time. DDPG is retained as a continuous-action policy-learning comparator, but the manuscript no longer relies on a strong claim that the block morphology problem is a natural MDP.

We also aligned this wording with the reward and episode definitions in Section 2.4. The added references support the general RL background, while the method text explains the specific implementation used here: a fixed surrogate, normalized target states, absolute descriptor actions, and a finite query horizon. This keeps the citation response useful without importing assumptions from physical control problems that are not part of this study.
}
\responseheading{Revisions made in the manuscript.}
\responsepara{%
We made three corresponding revisions. First, we added reinforcement-learning and DDPG references. Second, we rewrote the DDPG formulation around state, action, reward, episode length, actor, critic, and target networks. Third, we softened the wording so that the manuscript does not overstate the task as a natural MDP.
}
\responseheading{Relevant revised manuscript and supporting material.}
\begin{revisionbox}
\revisionloc{Main manuscript, Section 1.2 paragraph.}
Deep reinforcement learning provides another family of search methods, but its role in morphology design requires careful framing. Sequential decision problems are commonly formalized as Markov Decision Processes in reinforcement learning. However, the morphology task in this study is not treated as a physical-time design-control process. DDPG is used as a serialized static black-box search method on a guarded surrogate response surface.
\revisionseparator
\revisionloc{Main manuscript, Section 2.4 paragraph.}
The DDPG environment serializes repeated queries to a static guarded surrogate. At the beginning of each episode, a random descriptor query is evaluated to construct the initial normalized three-target state. At each subsequent step, the actor outputs an absolute 12-dimensional normalized descriptor action. This action is not an incremental modification of the preceding block. The action is denormalized to descriptor bounds, evaluated by the guarded surrogate, and converted to the next normalized target state and scalar reward. The 40-step horizon therefore functions as a fixed sequence of policy-training queries on a static surrogate response surface.
\revisionseparator
\revisionloc{Main manuscript, reward formula excerpt.}
\RewardEquations
\end{revisionbox}
\responseheading{Local references.}
\begin{responsereferences}
\referenceitem{Puterman, M. L. (1994). \textit{Markov Decision Processes: Discrete Stochastic Dynamic Programming}. Wiley.}
\referenceitem{Sutton, R. S., \& Barto, A. G. (2018). \textit{Reinforcement Learning: An Introduction} (2nd ed.). MIT Press.}
\referenceitem{Wang, X., Wang, S., et al. (2024). Deep reinforcement learning: A survey. \textit{IEEE Transactions on Neural Networks and Learning Systems}, 35(4), 5064--5078. \responsedoi{https://doi.org/10.1109/tnnls.2022.3207346}}
\referenceitem{Mnih, V., Kavukcuoglu, K., et al. (2015). Human-level control through deep reinforcement learning. \textit{Nature}, 518(7540), 529--533. \responsedoi{https://doi.org/10.1038/nature14236}}
\referenceitem{Lillicrap, T. P., Hunt, J. J., et al. (2015). Continuous control with deep reinforcement learning. \textit{arXiv preprint arXiv:1509.02971}. \responsedoi{https://doi.org/10.48550/arXiv.1509.02971}}
\end{responsereferences}

\comment{R1-11}
\begin{commentbox}
Figure 3: The manuscript does not include the explanation of the Figure 3. Moreover, the authors need to explain the symbols and equations in the Figure 3.
\end{commentbox}
\responseheading{Response.}
\checkresponse{%
We thank the reviewer for this comment. We agree that the previous manuscript did not explain Fig. 3 sufficiently. The figure contained the actor-critic architecture and update equations, but the surrounding text did not define the state, action, reward, target networks, temporal-difference target, critic loss, and actor update in enough detail.

We revised Section 2.4 and the Fig. 3 caption. The revised text now defines the DDPG state as the normalized three-target vector, the action as the absolute 12-dimensional normalized descriptor query, and the reward as \(R=1-d_w\). It also explains the online actor, online critic, target actor, target critic, replay buffer, TD target, critic loss, actor update, and soft target-network update shown in Fig. 3.

We also added the reward equations directly in the manuscript. The equations define the normalized target coordinate \(z_i\), the ideal vector \(\mathbf{u}=[0,1,1]^\mathsf{T}\), the weighted normalized distance \(d_w\), and the reward \(R=1-d_w\). This removes the previous reward-offset wording and makes the symbols in Fig. 3 consistent with the method text.

The revision also connects the equations to the diagram. Readers can now follow one DDPG update from a descriptor action, through guarded surrogate prediction, to the reward and TD target used by the critic. This makes Fig. 3 an explanation of the actual surrogate-query implementation rather than a generic actor-critic schematic.
}
\responseheading{Revisions made in the manuscript.}
\responsepara{%
We made three corresponding revisions. First, we expanded Section 2.4 to define the DDPG components. Second, we redrew Fig. 3 as a static surrogate-query architecture diagram. Third, we added the corrected normalization and reward equations.
}
\responseheading{Relevant revised manuscript and supporting material.}
\begin{revisionbox}
\revisionloc{Main manuscript, Section 2.4 paragraph defining state/action/episode.}
The DDPG environment serializes repeated queries to a static guarded surrogate. At the beginning of each episode, a random descriptor query is evaluated to construct the initial normalized three-target state. At each subsequent step, the actor outputs an absolute 12-dimensional normalized descriptor action. This action is not an incremental modification of the preceding block. The action is denormalized to descriptor bounds, evaluated by the guarded surrogate, and converted to the next normalized target state and scalar reward.
\revisionseparator
\revisionloc{Main manuscript, reward equations.}
\RewardEquations
\revisionseparator
\revisionloc{Main manuscript, Fig.~3 equation explanation.}
Figure 3 summarizes the actor-critic update used in this surrogate-query environment. The state \(s\) is the normalized three-target vector, the action \(a\) is the normalized 12-descriptor query, \(r\) is the scalar reward, \(s_{\mathrm{next}}\) is the next normalized target state returned by the guarded surrogate, and \(d\) is the episode-termination indicator. The online critic \(Q(s,a)\) estimates the action value, while the target actor and target critic provide the delayed target used in the temporal-difference update. The critic minimizes the mean-squared error between \(Q(s,a)\) and the TD target, while the actor is updated by maximizing the critic value of the actor-generated action.
\revisionseparator
\revisionloc{Main manuscript, Fig.~3.}
\compactfig{0.86}{fig3.pdf}
\figcaptionstyle{\textbf{Revised Fig.~3.} DDPG actor-critic architecture and target-network update structure used in the surrogate-query environment. The diagram identifies the normalized target state, 12-descriptor action, guarded surrogate transition, scalar reward, replay buffer, online actor and critic, target actor and critic, TD target, critic loss, actor update, and soft target-network update.}
\end{revisionbox}

\comment{R1-12}
\begin{commentbox}
Figure 4: For parity plots, the ticks and limits of the x- and y-axes should be same. All plots should be rescaled. Please specify unit in each axis.
\end{commentbox}
\responseheading{Response.}
\checkresponse{%
We thank the reviewer for this comment. We agree that the figure and table units needed clearer labelling. This point matters because unit ambiguity makes it hard to distinguish raw targets, normalized coordinates, and morphology descriptors.

We revised Figs. 4, 5, and 9, Tables 1, 3, and 6, and Supplementary captions and tables to standardize target units, morphology units, one-to-one reference lines, and figure captions. The revised material now gives the reader a more direct account of what was changed and how the revision should be interpreted.

The revised figures and tables label EUIt in kWh m$^{-2}$ y$^{-1}$, EG in $10^6$ kWh y$^{-1}$ where that scale is used, H in hours, and morphology quantities by their unit or type. This correction changes presentation only; the underlying numerical results are unchanged.

The response also distinguishes unit correction from result revision, which is important because the numbers are retained while the labels and comparison guides become easier to audit.
}
\responseheading{Revisions made in the manuscript.}
\responsepara{%
We made three corresponding revisions. First, we standardized target-unit labels in the relevant figures and tables. Second, we revised comparison plots to include one-to-one reference lines where appropriate. Third, we checked Supplementary Information captions and tables against the same convention.
}
\responseheading{Relevant revised manuscript and supporting material.}
\begin{revisionbox}
\revisionloc{Main manuscript, Fig.~4.}
\compactfig{0.86}{data_and_surrogate_validation.pdf}
\figcaptionstyle{\textbf{Revised Fig.~4.} Descriptor coverage and cross-validated surrogate predictions.}
\revisionseparator
\revisionloc{Main manuscript, Table 3 excerpt with units in target definitions.}
\SurrogateMetricsTable
\revisionseparator
\revisionloc{Main manuscript, Fig.~9.}
\compactfig{0.86}{physical_cross_model_stress_test.pdf}
\figcaptionstyle{\textbf{Revised Fig.~9.} Cross-model stress test for 18 direct generated feasible blocks.}
\end{revisionbox}

\comment{R1-13}
\begin{commentbox}
Figure 6: Please explain what the episode means in the plots.
\end{commentbox}
\responseheading{Response.}
\checkresponse{%
We thank the reviewer for this comment. We agree that the DDPG episode definition could be mistaken for physical time evolution. This point matters because in this static surrogate benchmark an episode is a query sequence rather than an annual or dynamic building-operation process.

We revised Sections 2.4 and 3.2, Figs. 2 and 6, and Supplementary Section S3.1 to define one episode as 40 surrogate-query steps and align the training-curve captions with that definition. The revised material now gives the reader a more direct account of what was changed and how the revision should be interpreted.

The revised explanation makes the reward at each step depend on the surrogate-predicted normalized target vector and preference weighting. The episode axis therefore describes training progress in descriptor-query space, not a physical calendar or dynamic urban system.

This explanation is included because the word episode is familiar from sequential control, but here it only denotes the fixed number of surrogate interactions used for training.
}
\responseheading{Revisions made in the manuscript.}
\responsepara{%
We made three corresponding revisions. First, we rewrote the episode definition in Section 2.4. Second, we revised Fig. 2 and Fig. 6 captions around surrogate-query steps. Third, we added supporting detail on seed-level DDPG behaviour in Supplementary Section S3.1.
}
\responseheading{Relevant revised manuscript and supporting material.}
\begin{revisionbox}
\revisionloc{Main manuscript, Section 2.4 paragraph.}
The 40-step horizon therefore functions as a fixed sequence of policy-training queries on a static surrogate response surface.
\revisionseparator
\revisionloc{Main manuscript, Section 3.2 paragraphs.}
Figure 6 describes how the DDPG agents interact with the guarded surrogate during training. Across the three preference scenarios, rewards increase rapidly during early training and then enter a high-reward region with visible seed-to-seed variability. Because each episode contains a fixed sequence of 40 surrogate queries, the trajectories represent policy-training behaviour on a static response surface. They should therefore be read as search trajectories over descriptor queries, not as the physical evolution of one block design.

The later portions of the curves remain informative. The agents can repeatedly reach favourable surrogate regions, yet the best-to-final gap indicates that late-stage regression remains present. This behaviour is consistent with the retained-output strategy used in Table 4, where one best scalarized candidate is retained from each seed.
\revisionseparator
\revisionloc{Main manuscript, Fig.~6.}
\compactfig{0.86}{ddpg_training_dynamics.pdf}
\figcaptionstyle{\textbf{Revised Fig.~6.} DDPG training trajectories across the three preference scenarios.}
\revisionseparator
\revisionloc{Supplementary Material, Section S3.1 paragraph.}
Each DDPG scenario uses 20 seeds, 600 episodes, and 40 surrogate interactions per episode. The seed-level summaries record best reward, final reward, plateau episode, and best-to-final regression. These quantities describe policy-training behaviour under serialized surrogate queries on a static surrogate response surface.
\end{revisionbox}

\comment{R1-14}
\begin{commentbox}
It would be better to conduct more generalized case studies, including different locations (different weather conditions in China) or different building types (an energy-efficient building, an old building, etc.).
\end{commentbox}
\responseheading{Response.}
\checkresponse{%
We thank the reviewer for this constructive suggestion. We agree that the original single-climate presentation was too narrow for a study intended to support urban energy design interpretation. In response, we added a limited cross-climate physical sensitivity analysis while keeping the interpretation conservative.

The added analysis uses the same four generated block morphologies and evaluates them under three additional Chinese climate files: Harbin, Beijing, and Guangzhou. These climates represent severe-cold, cold, and hot-summer warm-winter conditions. The baseline morphology set is held fixed, the surrogate is not retrained, and the block selection is not changed after observing the climate results. This design allows the analysis to isolate weather sensitivity for selected morphologies rather than mixing climate effects with a new optimizer run or a new morphology selection process.

The added calculations produced 12 additional case-weather combinations. The station-level mean changes relative to the baseline weather are substantial for EUIt, especially in Harbin, where mean EUIt increases by 55.681 kWh m\(^{-2}\) y\(^{-1}\). Beijing increases mean EUIt by 13.467 kWh m\(^{-2}\) y\(^{-1}\), while Guangzhou decreases it by 2.972 kWh m\(^{-2}\) y\(^{-1}\). EG ordering remains stable across all three additional climates, with Spearman \(\rho=1.0\). H ordering is less stable, with Spearman \(\rho=0.2\), \(0.4\), and \(0.2\) for Beijing, Guangzhou, and Harbin. EUIt ordering also becomes unstable under Harbin, where Spearman \(\rho=0.2\).

We did not add new building types in this revision because that would require a new building-stock definition, envelope assumptions, occupancy schedules, and physical-model calibration. We instead state this as a limitation and future-work direction. The revised manuscript therefore responds to the location and weather concern with additional calculations, while keeping the building-type expansion outside the current model scope.
}
\responseheading{Revisions made in the manuscript.}
\responsepara{%
We made four corresponding revisions. First, we added the cross-climate calculation protocol in the Methodology. Second, we added climate result interpretation in the Results. Third, we added weather metadata, station-level changes, case-level climate variation, and rank-stability results in Supplementary Information Section S6. Fourth, we revised the Discussion and limitations to state that additional building types, climate-aware surrogates, storage, EV loads, and district-level energy-system interactions remain future extensions.
}
\responseheading{Relevant revised manuscript and supporting material.}
\begin{revisionbox}
\revisionloc{Main manuscript, Section 2.6 paragraph.}
The baseline weather file used in the physical workflow serves as the reference condition for the cross-climate sensitivity analysis. The cross-climate sensitivity set then evaluates the same four block morphologies under three additional weather files representing Harbin, Beijing, and Guangzhou. The surrogate is not retrained, and the morphology set is not reselected after observing climate results.
\revisionseparator
\revisionloc{Main manuscript, Section 3.5 paragraph.}
The climate analysis in Fig. 10 evaluates four selected blocks under three additional weather files relative to the baseline weather. The demand response is strongly climate-dependent, especially under the severe-cold Harbin weather. EG ordering remains stable across the additional climates, suggesting that the rooftop-generation proxy is less sensitive to the weather substitution for this four-block set. H ordering is less stable, which is consistent with its dependence on solar geometry, obstruction, and weather-specific direct-sun conditions.
\revisionseparator
\revisionloc{Main manuscript, Fig.~10.}
\compactfig{0.86}{cross_climate_sensitivity.pdf}
\figcaptionstyle{\textbf{Revised Fig.~10.} Cross-climate sensitivity of four selected blocks relative to baseline weather.}
\revisionseparator
\revisionloc{Supplementary Material, Table S5 excerpt.}
\ClimateTable
\revisionseparator
\revisionloc{Supplementary Material, Section S6.3 paragraph.}
The rank-stability results indicate target-specific climate sensitivity. EG has stable ordering across the three additional climates, with Spearman $\rho$ equal to 1.0 in each case. H shows lower ordering stability, with Spearman $\rho$ values of 0.2 for Beijing, 0.4 for Guangzhou, and 0.2 for Harbin. EUIt ordering remains stable for Beijing and Guangzhou but changes under Harbin, where Spearman $\rho$ is 0.2. These results support a limited climate-sensitivity interpretation for the selected four-block set.
\end{revisionbox}

\comment{R1-15}
\begin{commentbox}
I think it would be great if the authors could write this paper for a broader audience.
\end{commentbox}
\responseheading{Response.}
\checkresponse{%
We thank the reviewer for this suggestion. We agree that the previous manuscript was written too narrowly for readers who are interested in urban energy design but may not follow every detail of DDPG, HV, IGD, or surrogate-response construction. In the revision, we rewrote the paper so that the energy-design problem appears before the optimizer details.

The revised Introduction is now organized into Motivation, Literature review, and Contributions and paper organization. The Motivation subsection starts from block-scale energy demand, rooftop PV potential, winter sunlight access, Chinese urban renewal, and climate-resilient planning. The Literature review separates urban energy modelling, surrogate-assisted optimization, Pareto search, and DRL operational-control literature. The Contributions subsection is limited to three points and explains the workflow in terms of generated morphologies, shared surrogate comparison, feasible projection, and limited physical/climate-sensitive calculations.

We also made several presentation changes for readability. The front matter now contains a Nomenclature block rather than only an abbreviation list. Table 1 defines morphology descriptors compactly, and Table 2 separates evaluation modes so that analytic response generation, surrogate prediction, feasible projection, physical stress testing, and climate sensitivity are not confused. Figure captions were expanded where needed, especially for Fig. 1, Fig. 3, Fig. 6, Fig. 9, and Fig. 10. The Discussion was rewritten to explain what the descriptor-space comparison supports and what cannot be inferred from the physical and climate-sensitive calculations.

These changes are intended to make the paper readable as an urban energy design study, not only as an optimization benchmark. At the same time, the manuscript remains cautious: it does not claim physical optimizer certification or broad cross-climate generalization.
}
\responseheading{Revisions made in the manuscript.}
\responsepara{%
We made five corresponding revisions. First, we reorganized the Introduction around motivation, literature review, and contributions. Second, we replaced the front-matter abbreviation list with Nomenclature. Third, we revised the method tables to reduce terminology ambiguity. Fourth, we expanded key figure captions and explanations. Fifth, we rewrote the Discussion and Conclusions to separate descriptor-space findings, feasible-projection interpretation, physical stress-test results, and climate sensitivity.
}
\responseheading{Relevant revised manuscript and supporting material.}
\begin{revisionbox}
\revisionloc{Main manuscript, Introduction motivation excerpt.}
Urban energy planning is increasingly shaped by the dual requirement of reducing building-sector emissions and maintaining reliable energy services. At the district and block scales, urban morphology controls floor area, envelope exposure, roof availability, obstruction, and solar access, which together affect both energy demand and distributed generation potential.
\revisionseparator
\revisionloc{Main manuscript, contribution paragraph.}
This study develops a surrogate-assisted comparison workflow for block-scale urban energy design. The contributions are threefold: the study establishes a morphology-to-performance workflow in which feasible residential block layouts are generated before 12 morphology descriptors are calculated; the study compares DDPG and NSGA-II under the same selected DNN surrogate and matched query budget; and the study connects surrogate-based optimization with repeated surrogate assessment, feasible-block projection, a 24-case physics-based cross-model stress test, and a four-block cross-climate sensitivity analysis.
\revisionseparator
\revisionloc{Main manuscript, Nomenclature excerpt.}
\begin{center}
\footnotesize
\begin{tabularx}{0.96\linewidth}{@{}p{0.14\linewidth}X p{0.13\linewidth}X@{}}
\toprule
\multicolumn{4}{@{}l}{\textit{Abbreviations}}\\
EUIt & Energy Use Intensity & EG & Energy Generation \\
H & winter sunlight hours & DNN & deep neural network \\
DDPG & Deep Deterministic Policy Gradient & NSGA-II & Non-dominated Sorting Genetic Algorithm II \\
CMA-ES & Covariance Matrix Adaptation Evolution Strategy & GHI & global horizontal irradiance \\
HV & Hypervolume & IGD & Inverted Generational Distance \\
nMAE & normalized mean absolute error & OSLI & open-space location index \\
\bottomrule
\end{tabularx}
\end{center}
\revisionseparator
\revisionloc{Main manuscript, Table 2 excerpt.}
\EvaluationModeTable
\revisionseparator
\revisionloc{Main manuscript, Discussion paragraph.}
These two sensitivity analyses do not convert the descriptor-space comparison into a physical ranking result. They instead define where the analytic response surface, generated feasible blocks, and selected physical workflows agree only partially, which is the boundary needed for early-stage design-support interpretation.
\end{revisionbox}

\reviewerpage{2}
\checkresponse{%
We thank Reviewer 2 for identifying the central issues around analytic approximation, physical-probe terminology, validation scope, sample generation, RL formulation, evaluation terminology, target formulas, and reproducibility. These comments shaped the largest part of the revision.
}

\comment{R2-1}
\begin{commentbox}
The manuscript states that the dataset is generated through an "analytic approximation" of the original EnergyPlus/Radiance workflow rather than through direct simulation. Please provide reasons for this workflow.
\end{commentbox}
\responseheading{Response.}
\checkresponse{%
We thank the reviewer for this comment. We agree that the reason for using an analytic response dataset needed to be explicit. This point matters because the analytic response surface is the controlled benchmark layer that makes optimizer comparison auditable before selected physical calculations are added.

We revised Sections 2.1--2.3, Table 2, and Supplementary Section S1.4 to define generated feasible morphologies, analytic target assignment, selected DNN surrogate training, and the boundary between analytic and physical assessment. The revised material now gives the reader a more direct account of what was changed and how the revision should be interpreted.

The supporting material documents the target functions, stochastic terms, clipping ranges, descriptor dependencies, nested sample regime, and common surrogate checkpoint used for DDPG and NSGA-II. The analytic response dataset is not presented as direct physical simulation of every candidate.

This also explains why the revision does not simply add more physical terminology. The analytic dataset is useful precisely because it is controlled and reproducible, while the selected physical cases are used later to test where that controlled benchmark becomes fragile.
}
\responseheading{Revisions made in the manuscript.}
\responsepara{%
We made three corresponding revisions. First, we rewrote Sections 2.1--2.3 to distinguish generated morphologies, analytic target assignment, and surrogate training. Second, we revised Table 2 to name analytic response generation as its own evaluation mode. Third, we added equations, stochastic terms, clipping ranges, and descriptor dependencies in Supplementary Section S1.4.
}
\responseheading{Relevant revised manuscript and supporting material.}
\begin{revisionbox}
\revisionloc{Main manuscript, Section 2.2 paragraphs.}
This module provides a controlled morphology-to-performance response surface for optimizer comparison. The nested 500/1000/1500/2000-sample pools are deterministic subsets of the same generated morphology pool under the recorded random seed and morphology-generation protocol. The physics-based case calculations are conducted subsequently on locked generated blocks to assess how the analytic target values relate to EnergyPlus- and Radiance-component outputs in selected cases.

The response directions follow the design objectives: lower EUIt is preferred, while higher EG and H are preferred. The active 2000-sample pool is the largest nested sample regime used for the selected DNN surrogate. The full equations, station-bias convention, stochastic terms, and clipping ranges are given in Supplementary Section S1. Table 2 distinguishes analytic response assignment, surrogate prediction, feasible projection, and the post-optimization sensitivity analyses.
\revisionseparator
\revisionloc{Main manuscript, Table 2 excerpt.}
\EvaluationModeTable
\revisionseparator
\revisionloc{Supplementary Material, Section S1.4 opening paragraph.}
The analytic response generator is the target-assignment module that maps a generated morphology descriptor vector, denoted $\mathbf{x}$, to three target responses, denoted $\mathbf{y}^{ana}=[\mathrm{EUIt},\mathrm{EG},H]^\mathsf{T}$. The mapping uses explicit formulas that combine density, openness, aspect-ratio, sky-view, orientation, and open-space-position terms to represent demand, rooftop-generation, and sunlight-access responses.
\end{revisionbox}

\comment{R2-2}
\begin{commentbox}
The study repeatedly refers to "physical probe" and "physical-stack probe," but the description suggests that the probe is still simplified and only partially aligned with the original metrics.
\end{commentbox}
\responseheading{Response.}
\checkresponse{%
We thank the reviewer for this comment. We agree that the physical-probe terminology was too strong. This point matters because selected physical calculations should test cross-model fragility rather than validate the entire surrogate or optimizer ranking.

We revised Table 2, Sections 2.6 and 3.5, Table 6, Fig. 9, and Supplementary Sections S5.3--S5.4 to replace physical-probe wording with limited physics-based cross-model stress test and rewrote the interpretation around selected locked cases. The revised material now gives the reader a more direct account of what was changed and how the revision should be interpreted.

The revised stress-test section compares EnergyPlus- and Radiance-component outputs against analytic target values for those cases. The stress test identifies weak agreement and ranking fragility; it does not support direct physical transfer of optimizer rankings.

The response therefore treats the terminology change as a claim-boundary correction. It explains why the calculation is still useful, but only as a cross-model stress test of selected cases rather than as broad validation.
}
\responseheading{Revisions made in the manuscript.}
\responsepara{%
We made three corresponding revisions. First, we replaced validation-style physical-probe wording with limited physics-based cross-model stress test. Second, we revised Table 2 and Section 2.6 to state the source, output, and boundary of the physical cases. Third, we rewrote the Results and Discussion around weak agreement rather than confirmation.
}
\responseheading{Relevant revised manuscript and supporting material.}
\begin{revisionbox}
\revisionloc{Main manuscript, Table 2 excerpt.}
\EvaluationModeTable
\revisionseparator
\revisionloc{Main manuscript, Section 2.6 paragraph.}
The limited physics-based cross-model stress test uses 24 locked cases. Eighteen direct feasible cases compare analytic response-generator values with physics-based outputs without optimizer projection, while six optimizer-linked cases decompose the path from descriptor-space candidate to nearest generated feasible block and then to physics-based result. The direct and optimizer-linked subsets answer different questions and are not merged into a single parity plot.
\revisionseparator
\revisionloc{Main manuscript, Section 3.5 paragraph.}
The physics-based cross-model stress test completed all locked cases and provides a case-level comparison between analytic target values and physics-based outputs. Figure 9 uses only the 18 direct feasible blocks, so the parity plots are not affected by optimizer projection. The three targets show limited metric and rank consistency in this case set.
\revisionseparator
\revisionloc{Supplementary Material, Section S5.4 paragraphs.}
The locked physical workflow did not produce annual roof-plane irradiance outputs for the direct stress-test cases. The unavailable roof-irradiance pathway is therefore kept separate from the completed EUIt, EG proxy, and H calculations.

EG is summarized with the GHI-based rooftop-PV proxy in the stress-test table. This proxy preserves a generation-related comparison across the locked cases, but it should be interpreted as a simplified rooftop-PV indicator rather than a full annual roof-irradiance result.
\end{revisionbox}

\comment{R2-3}
\begin{commentbox}
The validation evidence remains insufficient. Only a small-batch probe is used, and the physical results show noticeable deviations from surrogate predictions, especially for EUIt and H.
\end{commentbox}
\responseheading{Response.}
\checkresponse{%
We thank the reviewer for this comment. We agree that the previous validation wording was insufficient and too optimistic. This point matters because the physical calculations should be reported as weak agreement where the tested cases show that result.

We revised Sections 2.6, 3.5, and 4.2, Fig. 9, Table 6, and Supplementary Section S5 to report completed physical cases, direct feasible-case metrics, and the downgraded interpretation. The revised material now gives the reader a more direct account of what was changed and how the revision should be interpreted.

The added calculations show 18 direct feasible cases with EUIt MAE 3.788, nMAE 0.340, and Spearman rho -0.404; EG proxy MAE 0.441, nMAE 0.405, and rho -0.023; and H MAE 1.108, nMAE 0.601, and rho -0.043. These values support a stress-test conclusion about fragility, not physical equivalence or optimizer-ranking transfer.

This is why the response gives the numerical agreement values directly. The revision does not hide the weak correlations; it uses them to narrow the manuscript's conclusion and to prevent ranking-transfer language from reappearing.
}
\responseheading{Revisions made in the manuscript.}
\responsepara{%
We made four corresponding revisions. First, we added the completed-case count and direct feasible-case summary. Second, we reported MAE, nMAE, and Spearman rho values in Table 6 and Fig. 9. Third, we revised Supplementary Section S5 to give case-level summaries. Fourth, we rewrote the Discussion to state that physical agreement is weak.
}
\responseheading{Relevant revised manuscript and supporting material.}
\begin{revisionbox}
\revisionloc{Main manuscript, Section 3.5 paragraph.}
The physics-based cross-model stress test completed all locked cases and provides a case-level comparison between analytic target values and physics-based outputs. Figure 9 uses only the 18 direct feasible blocks, so the parity plots are not affected by optimizer projection. The three targets show limited metric and rank consistency in this case set, with the weakest ordering stability appearing in the sunlight-related and demand-related comparisons. This result does not diminish the value of the stress test; instead, it clarifies the degree to which the analytic response surface and the physics-based calculations differ for selected generated blocks.
\revisionseparator
\revisionloc{Main manuscript, Fig.~9.}
\compactfig{0.86}{physical_cross_model_stress_test.pdf}
\figcaptionstyle{\textbf{Revised Fig.~9.} Cross-model stress test for 18 direct generated feasible blocks.}
\revisionseparator
\revisionloc{Main manuscript and Supplementary Material, physical summary excerpt.}
\PhysicalSummaryTable
\revisionseparator
\revisionloc{Supplementary Material, Section S5.2 paragraph.}
The physics-based stress-test set contains 24 locked cases. The 18 direct feasible cases are used to compare analytic response-generator values with physics-based outputs without optimizer projection. The six optimizer-linked cases decompose the path from descriptor-space candidate to nearest generated block and then to physics-based result. These two subsets should not be merged into a single parity plot because they answer different questions.
\end{revisionbox}

\comment{R2-4}
\begin{commentbox}
The manuscript claims that the final benchmark uses the most accurate 2000-sample surrogate selected from a 500/1000/1500/2000-sample scale study. However, the sample-generation strategy is not sufficiently described. The authors should specify whether sampling was random, Latin hypercube, grid-based, constraint-based, or derived from feasible morphology generation. The distribution and coverage of the 12-dimensional design space should also be shown.
\end{commentbox}
\responseheading{Response.}
\checkresponse{%
We thank the reviewer for this comment. We agree that the random feasible morphology generation procedure needed more detail. This point matters because the 2000 samples should be understood as generated feasible morphologies represented by derived descriptors, not arbitrary points in an unconstrained descriptor cube.

We revised Section 2.1, Section 3.1, Table 1, Fig. 4, and Supplementary Sections S1.1--S1.5 to expand the generation protocol, descriptor coverage checks, 3 by 3 OSLI encoding, PCA summary, and duplicate checks. The revised material now gives the reader a more direct account of what was changed and how the revision should be interpreted.

The supporting material records prototype settings, descriptor distributions, six components for 95 percent descriptor variance, and absence of exact duplicate descriptor rows in the 2000-row canonical pool. This is why optimizer outputs are projected back to generated feasible blocks before morphology-level claims are made.

The response also connects generation to later projection. Because the optimizer searches descriptor space while the design object is a generated block, sample coverage and duplicate checks are necessary for interpreting feasible-morphology diversity.
}
\responseheading{Revisions made in the manuscript.}
\responsepara{%
We made four corresponding revisions. First, we expanded the morphology-generation description in Section 2.1. Second, we revised Table 1 to define descriptors and OSLI encoding. Third, we added Fig. 4 and supporting distribution material. Fourth, we added Supplementary Section S1 material on prototype settings, PCA coverage, descriptor dependencies, and duplicate checks.
}
\responseheading{Relevant revised manuscript and supporting material.}
\begin{revisionbox}
\revisionloc{Main manuscript, Section 2.1 paragraphs.}
The morphology generator first constructs feasible block layouts and only then calculates descriptor values. Each block is represented as a 240 m $\times$ 240 m square divided into a $3\times3$ grid of 80 m $\times$ 80 m land units. One cell is assigned as open space, and the remaining eight cells are populated with residential prototypes. For each generated block, the generator samples the open-space position, building prototype, number of floors, and orientation angle. The orientation angle is sampled within -45 deg to 45 deg.
\revisionseparator
\revisionloc{Main manuscript, Fig.~4.}
\compactfig{0.86}{data_and_surrogate_validation.pdf}
\figcaptionstyle{\textbf{Revised Fig.~4.} Descriptor coverage and cross-validated surrogate predictions.}
\revisionseparator
\revisionloc{Supplementary Material, Section S1.2 paragraph.}
For each generated block, the open-space index is sampled from 0 to 8, theta is sampled uniformly from -45 deg to 45 deg, and each populated cell receives a randomly sampled prototype and a floor count sampled within that prototype range. Descriptor values are calculated only after this geometry-generation step. This order keeps the generated feasible block as the primary object and treats the 12 surrogate inputs as derived morphology descriptors.
\revisionseparator
\revisionloc{Supplementary Material, Table S2 and Fig.~S1 excerpts.}
\begin{center}
\footnotesize
\begin{tabular}{lll}
\toprule
Check & Value & Interpretation \\
\midrule
Exact duplicate count & 0 & no exact duplicate descriptor rows \\
Nearest-neighbour distance median & 0.199 & sparse normalized descriptor space \\
PCA components for 95\% variance & 6 & structured descriptor space \\
\bottomrule
\end{tabular}
\end{center}
\compactfig{0.80}{A1_descriptor_distributions.pdf}
\end{revisionbox}

\comment{R2-5}
\begin{commentbox}
The RL formulation remains questionable because the problem is essentially a static black-box optimization problem rather than a naturally sequential decision-making process. The state transition is determined by the surrogate response to a morphology vector, and the sequential meaning of episodes and steps is unclear.
\end{commentbox}
\responseheading{Response.}
\checkresponse{%
We thank the reviewer for this comment. We agree that the DDPG static black-box search formulation needed to be transparent. This point matters because generic RL language can mislead readers if the static surrogate-query setting is not stated.

We revised Sections 1.2, 2.4, and 3.2, Figs. 2, 3, and 6, and Supplementary Section S3.1 to define state, action, episode, reward, actor, critic, replay buffer, and target networks for the static surrogate benchmark. The revised material now gives the reader a more direct account of what was changed and how the revision should be interpreted.

The revised text states that one episode is a sequence of 40 surrogate queries and that the reward is computed from the surrogate-predicted normalized target vector. DDPG is included as a policy-learning comparator under a matched query budget, not as a model of dynamic building operation.

This is why the response names each RL component instead of only citing DDPG. The component definitions show exactly where the formulation departs from physical time-series control and why the comparison remains a surrogate-query benchmark. It also explains why the episode count, query budget, and reward definition are reported together rather than treated as separate implementation details.
}
\responseheading{Revisions made in the manuscript.}
\responsepara{%
We made three corresponding revisions. First, we rewrote the DDPG method section around state, action, episode, reward, actor, critic, and target networks. Second, we revised Figs. 2, 3, and 6 to align with static surrogate-query interpretation. Third, we added Supplementary Section S3.1 for seed-level training behaviour.
}
\responseheading{Relevant revised manuscript and supporting material.}
\begin{revisionbox}
\revisionloc{Main manuscript, Section 2.4 paragraph.}
The DDPG environment serializes repeated queries to a static guarded surrogate. At the beginning of each episode, a random descriptor query is evaluated to construct the initial normalized three-target state. At each subsequent step, the actor outputs an absolute 12-dimensional normalized descriptor action. This action is not an incremental modification of the preceding block. The action is denormalized to descriptor bounds, evaluated by the guarded surrogate, and converted to the next normalized target state and scalar reward. The 40-step horizon therefore functions as a fixed sequence of policy-training queries on a static surrogate response surface.
\revisionseparator
\revisionloc{Main manuscript, Fig.~2.}
\compactfig{0.86}{fig2.pdf}
\figcaptionstyle{\textbf{Revised Fig.~2.} Serialized DDPG surrogate-query workflow. Each episode contains 40 absolute descriptor queries.}
\revisionseparator
\revisionloc{Main manuscript, reward equations.}
\RewardEquations
\revisionseparator
\revisionloc{Supplementary Material, Section S3.1 paragraph.}
Each DDPG scenario uses 20 seeds, 600 episodes, and 40 surrogate interactions per episode. The seed-level summaries record best reward, final reward, plateau episode, and best-to-final regression. These quantities describe policy-training behaviour under serialized surrogate queries on a static surrogate response surface.
\end{revisionbox}

\comment{R2-6}
\begin{commentbox}
The manuscript uses several terms that need clearer distinction, including "analytic response-generation workflow," "independent analytic reevaluation," "physical probe," "full physical stack," and "publication-grade physical closure." These terms are currently confusing. A table comparing these evaluation modes would improve transparency.
\end{commentbox}
\responseheading{Response.}
\checkresponse{%
We thank the reviewer for this comment. We agree that evaluation terminology was inconsistent. This point matters because readers need to know whether a result comes from analytic target assignment, surrogate prediction, projection, selected physical calculation, or climate sensitivity.

We revised Table 2 and Sections 2.2, 2.6, 3.5, and 4.2 to standardize the evaluation terms and assign each one a source, output, and interpretation boundary. The revised material now gives the reader a more direct account of what was changed and how the revision should be interpreted.

The revised text uses analytic response generation for EUIt/EG/H target assignment, surrogate prediction for the DNN layer, feasible projection for descriptor-to-block mapping, physical stress test for locked cases, and climate sensitivity for weather-file changes. This terminology prevents overclaiming by keeping supporting layers separate.

This terminology cleanup also supports the reviewer-response supporting blocks because each displayed excerpt now has a named source layer: analytic formula, surrogate result, projection diagnostic, physical stress test, or climate sensitivity.
}
\responseheading{Revisions made in the manuscript.}
\responsepara{%
We made three corresponding revisions. First, we rewrote Table 2 as an evaluation-mode table. Second, we aligned Sections 2.2, 2.6, 3.5, and 4.2 with the same terminology. Third, we removed ambiguous wording that could conflate analytic response assignment with physical simulation or broad validation.
}
\responseheading{Relevant revised manuscript and supporting material.}
\begin{revisionbox}
\revisionloc{Main manuscript, Table 2 excerpt.}
\EvaluationModeTable
\revisionseparator
\revisionloc{Main manuscript, Section 2.6 paragraph.}
Descriptor-space optimizer outputs are treated first as candidate descriptors. Each retained descriptor-space candidate is therefore mapped to the nearest generated feasible block in normalized descriptor space. This mapping quantifies projection distance, duplicate collapse, and changes in HV and IGD under the fixed normalized reference front. The projection step is a representation check that precedes morphology-level interpretation.
\revisionseparator
\revisionloc{Main manuscript, Discussion boundary paragraph.}
These two sensitivity analyses do not convert the descriptor-space comparison into a physical ranking result. They instead define where the analytic response surface, generated feasible blocks, and selected physical workflows agree only partially, which is the boundary needed for early-stage design-support interpretation.
\end{revisionbox}

\comment{R2-7}
\begin{commentbox}
The objective definitions require further clarification. EUIt, EG, and H are said to conceptually correspond to EnergyPlus, rooftop PV, and sunlight-hour calculations, but the actual values are generated analytically. The authors should provide clear formulas or workflows for calculating each target in the generated dataset.
\end{commentbox}
\responseheading{Response.}
\checkresponse{%
We thank the reviewer for this comment. We agree that the objective formulas and clipping ranges needed to be shown. This point matters because reproducibility and interpretation require the reader to see how EUIt, EG, and H are generated for the analytic dataset.

We revised Section 2.2, Table 2, and Supplementary Section S1.4 to add the analytic response-generation equations, stochastic terms, objective directions, and clipping ranges. The revised material now gives the reader a more direct account of what was changed and how the revision should be interpreted.

The revised Nomenclature now defines the orientation penalty, aspect-ratio average, openness and solar-access proxy terms, density penalty, response-generation noise, and final clipped outputs. The formulas define a controlled benchmark layer and are not claimed to be a full physical energy model.

The response also states the remaining boundary: exposing the formulas improves reproducibility of the analytic benchmark, but those equations are still target-assignment rules rather than a calibrated physical simulation model. This is why the response keeps both the equations and the clipping ranges in the point-by-point response package: the formulas define the target surface, while the ranges define the benchmark values seen by the surrogate and optimizers.
}
\responseheading{Revisions made in the manuscript.}
\responsepara{%
We made three corresponding revisions. First, we added the EUIt, EG, and H analytic formulas to Supplementary Section S1.4. Second, we added the clipping ranges and stochastic-term convention. Third, we revised Section 2.2 and Table 2 so that the formulas are identified as analytic targets used for surrogate training and optimizer comparison.
}
\responseheading{Relevant revised manuscript and supporting material.}
\begin{revisionbox}
\revisionloc{Main manuscript, Section 2.2 paragraph.}
The response directions follow the design objectives: lower EUIt is preferred, while higher EG and H are preferred. The active 2000-sample pool is the largest nested sample regime used for the selected DNN surrogate. The full equations, station-bias convention, stochastic terms, and clipping ranges are given in Supplementary Section S1. Table 2 distinguishes analytic response assignment, surrogate prediction, feasible projection, and the post-optimization sensitivity analyses.
\revisionseparator
\revisionloc{Supplementary Material, Section S1.4 equations.}
\AnalyticEquations
\end{revisionbox}

\comment{R2-8}
\begin{commentbox}
The paper should provide code, data, or at least detailed reproducibility information.
\end{commentbox}
\responseheading{Response.}
\checkresponse{%
We thank the reviewer for this comment. We agree that the reproducibility statement needed to be clearer about public and excluded material. This point matters because the response package should identify what can be released without exposing credentials, local logs, or restricted raw inputs.

We revised the Data and code availability statement, README, and response manifest to rewrite the availability wording around a sanitized public GitHub release. The revised material now gives the reader a more direct account of what was changed and how the revision should be interpreted.

The revised text states that scripts and processed tables for the analytic response dataset, surrogate assessment, optimizer comparison, feasible projection, physical stress-test summaries, and cross-climate sensitivity summaries will be made available. Raw credentials, private logs, machine-specific paths, and non-redistributable inputs remain excluded.

This wording is intentionally practical. It tells the reviewer what can be audited from the public release while avoiding any promise to publish credentials, private paths, or raw material that cannot be redistributed.
}
\responseheading{Revisions made in the manuscript.}
\responsepara{%
We made three corresponding revisions. First, we rewrote the Data and code availability statement. Second, we aligned the response package README and manifest with the files used for this rebuttal. Third, we preserved the exclusion boundary for credentials, local logs, private paths, and restricted raw inputs.
}
\responseheading{Relevant revised manuscript and supporting material.}
\begin{revisionbox}
\revisionloc{Main manuscript, Data and code availability.}
The code, processed datasets, and result tables supporting the numerical analyses will be made available through the project repository at \href{https://github.com/haolinwang1313/DRL-Driven-Optimization}{https://github.com/haolinwang1313/DRL-Driven-Optimization}. The public release will include the scripts required to regenerate the analytic response dataset, surrogate assessment summaries, optimizer comparison tables, feasible-projection results, physics-based stress-test summaries, and cross-climate sensitivity results.
\end{revisionbox}

\reviewerpage{3}
\checkresponse{%
We thank Reviewer 3 for the comments on workflow-figure spelling, utility interpretation, reward formulation, unit consistency, and analytic/physical wording. These comments helped align the figures, equations, and claims with the revised material.
}

\comment{R3-1}
\begin{commentbox}
Correct spelling and wording errors in the workflow figures. For example, Fig. 1 still contains visible typos such as "loactions," "Acotr netword," and "Critic netword," which should be revised to "locations," "Actor network," and "Critic network."
\end{commentbox}
\responseheading{Response.}
\checkresponse{%
We thank the reviewer for this comment. We agree that the figure spelling and label errors needed correction. This point matters because workflow and architecture figures define the method structure, so visible errors reduce clarity.

We revised Fig. 1, Fig. 3, and their captions to correct the typo and align figure labels with the revised methodology and DDPG terminology. The revised material now gives the reader a more direct account of what was changed and how the revision should be interpreted.

Fig. 1 now uses consistent workflow-module labels, while Fig. 3 uses standardized labels for actor, critic, replay buffer, target networks, state, action, and reward path. This is a presentation correction and does not change the method or results.
}
\responseheading{Revisions made in the manuscript.}
\responsepara{%
We made two corresponding revisions. First, we corrected the typo and terminology issues in Fig. 1 and its caption. Second, we corrected Fig. 3 and its caption so that the DDPG labels match the revised method description.
}
\responseheading{Relevant revised manuscript and supporting material.}
\begin{revisionbox}
\revisionloc{Main manuscript, Fig.~1.}
\compactfig{0.86}{fig1.pdf}
\figcaptionstyle{\textbf{Revised Fig.~1.} Workflow linking feasible morphology generation, analytic response assignment, surrogate-assisted optimization, feasible-block projection, and sensitivity analysis.}
\revisionseparator
\revisionloc{Main manuscript, Fig.~3.}
\compactfig{0.86}{fig3.pdf}
\figcaptionstyle{\textbf{Revised Fig.~3.} DDPG actor-critic architecture and target-network update structure used in the surrogate-query environment.}
\end{revisionbox}

\comment{R3-2}
\begin{commentbox}
Verify the consistency between Fig. 9(d) and the stated benchmark conclusion. The text states that NSGA-II remains stronger than all DDPG scenarios, but the utility panel should be carefully checked to ensure that the plotted curves, legend, and utility direction support this conclusion.
\end{commentbox}
\responseheading{Response.}
\checkresponse{%
We thank the reviewer for this comment. We agree that the earlier utility interpretation could be read as a direct optimizer-ranking result. This point matters because benchmark fairness requires separating scalarized utility, equal-size HV/IGD, retained outputs, and feasible projection.

We revised Sections 3.3 and 3.4, Table 5, Figs. 7 and 8, and Supplementary Sections S3 and S5 to rewrite the comparison around fixed-reference descriptor-space metrics, equal-size diagnostics, and projection. The revised material now gives the reader a more direct account of what was changed and how the revision should be interpreted.

The added calculations show NSGA-II strong in descriptor-space HV and IGD under the selected surrogate benchmark while DDPG remains preference dependent and projection changes feasible-block interpretation. Projection and physical stress-test results prevent descriptor-space ranking from being read as direct physical performance ranking.

The response therefore avoids ranking language based on one utility view. It explains that each metric has its own contract and that the projection stage changes how descriptor-space performance should be interpreted.
}
\responseheading{Revisions made in the manuscript.}
\responsepara{%
We made three corresponding revisions. First, we rewrote Sections 3.3 and 3.4 to separate scalarized candidate selection, descriptor-space HV and IGD, equal-size diagnostics, and feasible projection. Second, we revised Table 5 and Figs. 7--8 around those criteria. Third, we removed the previous utility-only interpretation from the main claim.
}
\responseheading{Relevant revised manuscript and supporting material.}
\begin{revisionbox}
\revisionloc{Main manuscript, Section 3.3 paragraphs.}
Table 5 summarizes descriptor-space performance under the fixed normalized reference front, while Fig. 7 separates fixed-domain utility from equal-size archive metrics. NSGA-II is close to the HV ceiling under the fixed reference and remains the strongest method in the equal-size HV and IGD comparison. Among the three DDPG settings, the Generation-focused scenario gives the strongest descriptor-space archive metrics, followed by the Energy Saving and Balanced settings.

The comparison should be interpreted through the retained-output structure. DDPG contributes one best scalarized candidate from each seed, whereas NSGA-II contributes final populations from repeated runs. The equal-size summary reduces, but does not eliminate, this structural difference. Therefore, the descriptor-space result supports a method-comparison statement under the selected surrogate and fixed reference; it does not by itself describe feasible-block diversity.
\revisionseparator
\revisionloc{Main manuscript, Fig.~7.}
\compactfig{0.86}{benchmark_fairness.pdf}
\figcaptionstyle{\textbf{Revised Fig.~7.} Descriptor-space DDPG and NSGA-II comparison under fixed-domain utility and equal-size retained outputs.}
\revisionseparator
\revisionloc{Main manuscript, Table 5 excerpt.}
\BenchmarkTable
\revisionseparator
\revisionloc{Main manuscript, Section 3.4 paragraph.}
The projection analysis shows that descriptor-space coverage and feasible-morphology diversity are distinct. NSGA-II retains 2000 descriptor candidates, yet nearest-neighbour projection maps them to 51 unique generated feasible blocks.
\end{revisionbox}

\comment{R3-3}
\begin{commentbox}
Recheck the mathematical form of the weighted-distance reward in Eq. (5). The current expression appears to apply weights directly to the normalized objective values; if the intended formulation is a standard weighted Euclidean distance, the equation should be revised accordingly.
\end{commentbox}
\responseheading{Response.}
\checkresponse{%
We thank the reviewer for this comment. We agree that the reward formula needed correction. This point matters because the previous expression could be misread as an artificial large-offset reward.

We revised Section 2.4 and Fig. 3 to define the reward directly from normalized weighted distance. The revised material now gives the reader a more direct account of what was changed and how the revision should be interpreted.

The revised equations normalize targets to $z_i$, uses the ideal vector $[0,1,1]^\mathsf{T}$, computes $d_w$ with the preference vector, and sets $R=1-d_w$ with no large offset. This aligns the manuscript with the static surrogate-query implementation and makes the reward scale interpretable.

The response also explains why the correction matters: a bounded distance-based reward is interpretable across preference settings, while the removed offset could obscure the actual optimization signal.
}
\responseheading{Revisions made in the manuscript.}
\responsepara{%
We made three corresponding revisions. First, we corrected the reward equations in Section 2.4. Second, we revised Fig. 3 so that the reward path follows normalization, ideal-vector comparison, weighted distance, and $R=1-d_w$. Third, we removed the confusing large-offset notation from the response and manuscript text.
}
\responseheading{Relevant revised manuscript and supporting material.}
\begin{revisionbox}
\revisionloc{Main manuscript, Section 2.4 reward equations.}
\RewardEquations
\revisionseparator
\revisionloc{Main manuscript, preference-weight sentence.}
The preference weights follow the target order [EUIt, EG, H]: Balanced uses $\mathbf{w}=[1/3,1/3,1/3]^\mathsf{T}$, Energy Saving uses $\mathbf{w}=[0.6,0.2,0.2]^\mathsf{T}$, and Energy Generation uses $\mathbf{w}=[0.2,0.6,0.2]^\mathsf{T}$.
\revisionseparator
\revisionloc{Main manuscript, reward boundary sentence.}
This reward is separate from the fixed-domain post-hoc utility used later for benchmark comparison. The fixed-domain utility is used for post-hoc comparison, while $R$ is the scalar signal used during serialized DDPG training.
\end{revisionbox}

\comment{R3-4}
\begin{commentbox}
Further standardize units in figure axes and table headings. Terms such as EUIt, EG, H, residuals, and morphology factors should consistently include units where applicable, for example "EUIt (kWh/m²/y)," "EG (10⁶ kWh/y)," and "H (h)."
\end{commentbox}
\responseheading{Response.}
\checkresponse{%
We thank the reviewer for this comment. We agree that target and morphology units were not consistently presented. This point matters because unit inconsistency makes it harder to audit results and compare figures.

We revised Tables 1, 3, 5, and 6, Figs. 4--10, and Supplementary tables and captions to standardize target units, morphology units, figure axes, and table headings. The revised material now gives the reader a more direct account of what was changed and how the revision should be interpreted.

The revised Abstract now reports EUIt as kWh m$^{-2}$ y$^{-1}$, EG using the stated $10^6$ kWh y$^{-1}$ scale where applicable, H in hours, and morphology quantities by unit or type. The numerical results are unchanged; the revision improves auditability.
}
\responseheading{Revisions made in the manuscript.}
\responsepara{%
We made two corresponding revisions. First, we standardized target and morphology units in the main figures and tables. Second, we aligned the Supplementary Information table and caption labels with the same unit convention.
}
\responseheading{Relevant revised manuscript and supporting material.}
\begin{revisionbox}
\revisionloc{Main manuscript, Table 1 descriptor units excerpt.}
\CompactDescriptorTable
\revisionseparator
\revisionloc{Main manuscript, Fig.~4 and Fig.~9 unit-bearing excerpts.}
\compactfig{0.80}{data_and_surrogate_validation.pdf}
\compactfig{0.80}{physical_cross_model_stress_test.pdf}
\figcaptionstyle{\textbf{Revised figures.} Target units are attached to axes and table headings where applicable.}
\revisionseparator
\revisionloc{Supplementary Material, Table S5 note.}
$\Delta$EUIt is in kWh m$^{-2}$ y$^{-1}$, $\Delta$EG is in $10^6$ kWh y$^{-1}$, and $\Delta$H is in h. Positive values denote increases relative to the baseline weather file.
\end{revisionbox}

\comment{R3-5}
\begin{commentbox}
Make the wording about simulation, analytic approximation, and physical probing fully consistent. Since the manuscript states that the dataset is based on an analytic approximation rather than direct full-stack EnergyPlus/Radiance simulations, all related phrases should avoid implying that the final benchmark was performed using direct physical simulation.
\end{commentbox}
\responseheading{Response.}
\checkresponse{%
We thank the reviewer for this comment. We agree that analytic, simulation, and physical wording needed consistent use. This point matters because the source layer changes the interpretation of every optimizer and validation statement.

We revised the Abstract, Sections 1--4, Table 2, Table 6, and Supplementary Sections S1 and S5 to standardize analytic response generation, surrogate prediction, physical stress test, and climate sensitivity terminology. The revised material now gives the reader a more direct account of what was changed and how the revision should be interpreted.

The revised text makes clear that analytic response generation assigns targets to the generated pool, the surrogate predicts those targets during optimization, and selected locked blocks are examined through physical calculations and weather-file changes. This distinction supports the revised boundary: surrogate rankings are meaningful within the selected benchmark, while physical stress-test results show weak agreement and unsupported ranking transfer.
}
\responseheading{Revisions made in the manuscript.}
\responsepara{%
We made three corresponding revisions. First, we standardized analytic response generation, surrogate prediction, physical stress test, and climate sensitivity terminology. Second, we revised Table 2 and Table 6 to carry the same source-layer distinction. Third, we removed wording that could imply that the 2000 analytic samples were direct physical simulations.
}
\responseheading{Relevant revised manuscript and supporting material.}
\begin{revisionbox}
\revisionloc{Main manuscript, Abstract excerpt.}
Feasible block layouts are generated first and converted into 12 morphology descriptors. A morphology-to-performance response surface then assigns three targets to each sample: Energy Use Intensity, Energy Generation, and winter sunlight hours.
\revisionseparator
\revisionloc{Main manuscript, Table 2 excerpt.}
\EvaluationModeTable
\revisionseparator
\revisionloc{Main manuscript, Discussion boundary paragraph.}
These two sensitivity analyses do not convert the descriptor-space comparison into a physical ranking result. They instead define where the analytic response surface, generated feasible blocks, and selected physical workflows agree only partially, which is the boundary needed for early-stage design-support interpretation.
\revisionseparator
\revisionloc{Supplementary Material, Section S5.4 paragraph.}
EG is summarized with the GHI-based rooftop-PV proxy in the stress-test table. This proxy preserves a generation-related comparison across the locked cases, but it should be interpreted as a simplified rooftop-PV indicator rather than a full annual roof-irradiance result.
\end{revisionbox}

\reviewerpage{4}
\checkresponse{%
We thank Reviewer 4 for comments on study necessity, climate-zone effects, energy-system expansion, best-scenario definition, and comparison with existing urban energy planning approaches. The revised manuscript strengthens the motivation while preserving the current model scope.
}

\comment{R4-1}
\begin{commentbox}
There are plenty of studies related to block-scale urban energy planning, the necessity of this study needs to be further highlighted.
\end{commentbox}
\responseheading{Response.}
\checkresponse{%
We thank the reviewer for this comment. We agree that the study necessity needed a clearer explanation. This point matters because block-scale planning decisions require early screening of urban form, energy demand, rooftop PV potential, and winter solar access before detailed system design is available.

We revised Sections 1.1--1.3 and the Discussion to rewrite the motivation, contribution, and positioning around representation-aware design support. The revised material now gives the reader a more direct account of what was changed and how the revision should be interpreted.

The added calculations show why optimizer rankings depend on descriptor coverage, retained-output contracts, feasible-block projection, surrogate checkpoint, and evaluation mode. The contribution is a design-support and benchmarking contribution, not a universal superiority claim for DDPG or NSGA-II.

The response now ties necessity to the failure mode that motivated the revision: optimizer rankings can look precise while still depending on representation, surrogate reliability, and the projection from descriptors back to feasible blocks. This is the reason the response emphasizes benchmark fragility as a research need rather than treating optimizer comparison as a routine leaderboard exercise.
}
\responseheading{Revisions made in the manuscript.}
\responsepara{%
We made three corresponding revisions. First, we rewrote the motivation and contribution paragraphs in the Introduction. Second, we added literature positioning for block-scale morphology, surrogate-assisted optimization, and policy-learning comparison. Third, we revised the Discussion so that study necessity is tied to representation-aware design support.
}
\responseheading{Relevant revised manuscript and supporting material.}
\begin{revisionbox}
\revisionloc{Main manuscript, Motivation paragraphs.}
Block morphology influences energy performance through coupled physical mechanisms. Density and height alter heating and cooling demand by changing envelope exposure, mutual shading, and microclimatic context. Solar access and daylight availability depend on roof exposure, sky-view conditions, block orientation, and open-space arrangement. Recent large-scale studies further show that urban spatial structure can shape building energy demand across broad urban areas and under extreme climate events, which makes morphology-sensitive design relevant beyond isolated building cases.

The computational challenge is that morphology-sensitive energy analysis becomes expensive when thousands of candidate blocks, multiple optimizer seeds, and several preference settings must be evaluated. Simulation-based urban form generation and building-performance optimization provide a strong basis for this task, but repeated full physical simulation remains costly in search loops.
\revisionseparator
\revisionloc{Main manuscript, gap paragraph.}
The remaining gap is methodological and representation-oriented. Urban energy optimization has established the value of morphology-aware design, surrogate modelling has reduced the cost of repeated evaluation, and reinforcement learning has proved useful in several energy-operation settings. What remains less clear is how policy learning and Pareto search should be compared when both operate on a shared surrogate response surface, when the surrogate inputs are descriptors derived from generated blocks, and when continuous outputs require projection before they can be interpreted as feasible morphologies.
\end{revisionbox}
\responseheading{Local references.}
\begin{responsereferences}
\referenceitem{Perera, A. T. D., Javanroodi, K., et al. (2023). Challenges resulting from urban density and climate change for the EU energy transition. \textit{Nature Energy}, 8(4), 397--412. \responsedoi{https://doi.org/10.1038/s41560-023-01232-9}}
\referenceitem{Shi, Z., Fonseca, J. A., \& Schlueter, A. (2017). A review of simulation-based urban form generation and optimization for energy-driven urban design. \textit{Building and Environment}, 121, 119--129. \responsedoi{https://doi.org/10.1016/j.buildenv.2017.05.006}}
\referenceitem{Westermann, P., \& Evins, R. (2019). Surrogate modelling for sustainable building design: A review. \textit{Energy and Buildings}, 198, 170--186. \responsedoi{https://doi.org/10.1016/j.enbuild.2019.05.057}}
\end{responsereferences}

\comment{R4-2}
\begin{commentbox}
The selected location for modelling is in a specific climate zone. How this specific climate condition may affect the modelling and how the findings could be translated into other climate zones? This is not explained clearly.
\end{commentbox}
\responseheading{Response.}
\checkresponse{%
We thank the reviewer for this comment. We agree that the climate-zone limitation needed direct treatment. This point matters because selected climate sensitivity is useful only if the manuscript states what it covers and what it does not cover.

We revised Sections 2.6, 3.5, and 4.2, Fig. 10, Table 6, and Supplementary Section S6 to add a four-block by three-climate sensitivity calculation and report weather metadata. The revised material now gives the reader a more direct account of what was changed and how the revision should be interpreted.

The added calculations show stable EG ordering in the tested four-block set but climate-dependent EUIt and H responses, with the severe-cold case producing the largest EUIt increase. The surrogate is not retrained by climate, and the analysis does not cover all climate zones or building types.

The response also states why the added analysis remains limited. It gives a climate sensitivity signal for selected locked blocks, but it does not claim that the surrogate or optimizer ranking generalizes across climate zones. The added wording also prevents climate sensitivity from being read as a new universal climate-transfer claim.
}
\responseheading{Revisions made in the manuscript.}
\responsepara{%
We made three corresponding revisions. First, we added the four-block by three-climate sensitivity design. Second, we added Fig. 10, Table 6, Supplementary Table S5, and weather metadata. Third, we revised the Discussion to state that broader climate-zone coverage and climate-aware surrogate training remain future work.
}
\responseheading{Relevant revised manuscript and supporting material.}
\begin{revisionbox}
\revisionloc{Main manuscript, Section 2.6 paragraph.}
The baseline weather file used in the physical workflow serves as the reference condition for the cross-climate sensitivity analysis. The cross-climate sensitivity set then evaluates the same four block morphologies under three additional weather files representing Harbin, Beijing, and Guangzhou. The surrogate is not retrained, and the morphology set is not reselected after observing climate results.
\revisionseparator
\revisionloc{Main manuscript, Fig.~10.}
\compactfig{0.86}{cross_climate_sensitivity.pdf}
\figcaptionstyle{\textbf{Revised Fig.~10.} Cross-climate sensitivity of four selected blocks relative to baseline weather.}
\revisionseparator
\revisionloc{Supplementary Material, Section S6.1 paragraph.}
The baseline weather file corresponds to the station used for the case-study physical workflow. The cross-climate sensitivity analysis evaluates the same four generated feasible blocks under Harbin, Beijing, and Guangzhou weather. The climate buckets are severe cold for Harbin, cold for Beijing, and hot-summer warm-winter for Guangzhou. The TMY period is 2009--2023 for these weather files.
\revisionseparator
\revisionloc{Supplementary Material, Table S5 excerpt.}
\ClimateTable
\revisionseparator
\revisionloc{Supplementary Material, Section S6.3 paragraph.}
EG has stable ordering across the three additional climates, with Spearman $\rho$ equal to 1.0 in each case. H shows lower ordering stability, with Spearman $\rho$ values of 0.2 for Beijing, 0.4 for Guangzhou, and 0.2 for Harbin. EUIt ordering remains stable for Beijing and Guangzhou but changes under Harbin, where Spearman $\rho$ is 0.2.
\end{revisionbox}

\comment{R4-3}
\begin{commentbox}
How the different sources of energy potentially affect the urban energy planning? This is specifically the case while China has rapidly developed the EV market. Can the distrbuted energy systems affect the urban energy planning?
\end{commentbox}
\responseheading{Response.}
\checkresponse{%
We thank the reviewer for this comment. We agree that EV charging, storage, and distributed energy are important but outside the current objective system. This point matters because adding them as if they were modeled would create an unsupported claim.

We revised Section 4.3 and the objective-scope discussion to state the current objective scope and add these systems as future extensions. The revised material now gives the reader a more direct account of what was changed and how the revision should be interpreted.

The current workflow covers feasible residential block morphologies, analytic EUIt/EG/H target assignment, surrogate comparison, projection, and selected physical and climate-sensitive cases. Storage dispatch, EV charging demand, district-level grid exchange, tariff response, and distributed-energy operation are not modeled in the current revision.

This handling keeps the revision honest: the manuscript acknowledges these energy-system interactions as important, but it does not retrofit them into objectives or calculations that were not part of the current study.
}
\responseheading{Revisions made in the manuscript.}
\responsepara{%
We made two corresponding revisions. First, we revised Section 4.3 to state that storage, EV charging, distributed-energy operation, and district grid exchange are outside the current model. Second, we added these systems to future work as extensions that require new objectives and calculations.
}
\responseheading{Relevant revised manuscript and supporting material.}
\begin{revisionbox}
\revisionloc{Main manuscript, Section 4.3 paragraph.}
Several limitations remain. The optimizer comparison is conducted on an analytic response-generation dataset, while direct physical calculations are available for the locked stress-test cases. The 12 surrogate inputs are morphology-derived descriptors rather than independent unconstrained design variables. The DDPG horizon is best understood as a policy-training query sequence on a static surrogate. The cross-model stress test indicates limited rank consistency in the tested case set, and the climate analysis covers four blocks and three additional climates. Future work should connect optimization to a feasible morphology decoder, expand the physical simulation set, incorporate storage, EV charging, and distributed-energy interactions, and train climate-aware surrogates across multiple regions.
\revisionseparator
\revisionloc{Main manuscript, objective scope excerpt.}
The workflow generates feasible residential block morphologies, converts them into 12 morphology descriptors, assigns EUIt, EG, and H through an explicit morphology-to-performance response surface, and trains a selected DNN surrogate from the 2000-sample regime of a nested scale study.
\end{revisionbox}

\comment{R4-4}
\begin{commentbox}
It is not very clear what is considered as "best scenario" for optimization process.
\end{commentbox}
\responseheading{Response.}
\checkresponse{%
We thank the reviewer for this comment. We agree that best scenario was ambiguous. This point matters because different criteria answer different questions and should not be collapsed into one label.

We revised Sections 2.5, 3.3, and 3.4, Tables 4--5, and Figs. 7--8 to separate best scalarized DDPG candidate, retained-output contracts, fixed-reference descriptor-space HV and IGD, equal-size comparison, feasible-block projection, and selected physical-case interpretation. The revised material now gives the reader a more direct account of what was changed and how the revision should be interpreted.

The added calculations show that NSGA-II is strong in descriptor-space HV and IGD under the fixed benchmark reference, while DDPG outcomes are preference dependent and projection changes morphology-level interpretation. No single statement from these criteria defines a universal best physical scenario.

The response also clarifies the practical consequence for readers: a candidate can be strong under one scalarized or archive metric and still require projection and selected physical-case checks before design interpretation.
}
\responseheading{Revisions made in the manuscript.}
\responsepara{%
We made three corresponding revisions. First, we rewrote Sections 2.5, 3.3, and 3.4 to define retained-output contracts and comparison criteria. Second, we revised Tables 4--5 and Figs. 7--8 to separate scalarized candidates, archive metrics, and projection. Third, we removed ambiguous best-scenario wording.
}
\responseheading{Relevant revised manuscript and supporting material.}
\begin{revisionbox}
\revisionloc{Main manuscript, Table 4 excerpt.}
\begin{center}
\footnotesize
\begin{tabular}{lrrl}
\toprule
Method & Runs & Queries/run & Retained output \\
\midrule
DDPG & 20 seeds & 24,000 & best scalarized point \\
NSGA-II & 20 runs & 24,000 & final populations \\
CMA-ES & 20 runs & 24,000 & supplementary archive \\
RandomSearch & 20 runs & 24,000 & supplementary archive \\
\bottomrule
\end{tabular}
\end{center}
\revisionseparator
\revisionloc{Main manuscript, Table 5 excerpt.}
\BenchmarkTable
\revisionseparator
\revisionloc{Main manuscript, Section 3.4 paragraph.}
Projection distance adds a second layer of interpretation. The DDPG scenarios retain fewer descriptor candidates than NSGA-II, but their projected candidates are generally farther from the generated feasible blocks. In contrast, NSGA-II shows stronger descriptor-space coverage and shorter projection distances, yet its candidates still collapse to a much smaller feasible-block set.
\revisionseparator
\revisionloc{Main manuscript, Fig.~8.}
\compactfig{0.86}{feasible_projection.pdf}
\figcaptionstyle{\textbf{Revised Fig.~8.} Descriptor-space candidate compression and projection distance after nearest-feasible-block mapping.}
\end{revisionbox}

\comment{R4-5}
\begin{commentbox}
More discussions are required to compare the proposed urban energy planning framework with existing approaches to validate its advantages.
\end{commentbox}
\responseheading{Response.}
\checkresponse{%
We thank the reviewer for this comment. We agree that the manuscript needed a clearer comparison with existing approaches. This point matters because the comparison should explain the workflow's contribution without overstating its advantage.

We revised Sections 1.2, 4.1, and 4.3 to expand the comparison with simulation-based optimization, surrogate-assisted optimization, NSGA-II Pareto archives, DRL control studies, and neighbouring derivative-free methods. The revised material now gives the reader a more direct account of what was changed and how the revision should be interpreted.

The revised text states the advantage narrowly as transparent separation of fast descriptor-space screening, feasible-block projection, and selected sensitivity analysis. The current results support representation-aware early-stage design support, not comprehensive superiority over existing approaches.

The response therefore compares mechanisms rather than making a broad advantage claim. The stated contribution is that the workflow exposes representation and evaluation boundaries that can be hidden in conventional optimizer comparisons. This keeps the comparison useful for reviewers who want to know what the framework adds while preserving the limits imposed by the current results. It also names the remaining work needed for any stronger comparison.
}
\responseheading{Revisions made in the manuscript.}
\responsepara{%
We made three corresponding revisions. First, we expanded the literature review to position the workflow relative to existing optimization and DRL studies. Second, we revised Section 4.1 to explain the representation-aware advantage. Third, we revised Section 4.3 to keep the comparison bounded by the current results.
}
\responseheading{Relevant revised manuscript and supporting material.}
\begin{revisionbox}
\revisionloc{Main manuscript, literature review paragraph.}
Parametric design studies show that geometry generation can support energy- and daylight-oriented design exploration. Research on parametric daylighting and low-energy design demonstrates that early-stage geometry and envelope choices can be formalized as searchable design spaces. Evolutionary multi-objective optimization remains widely used because it provides explicit trade-off archives for conflicting objectives. NSGA-II and related methods have been applied to envelope design, photovoltaic-oriented inverse urban design, urban block optimization, and energy-efficient shading design.
\revisionseparator
\revisionloc{Main manuscript, DRL comparison paragraph.}
In the energy domain, DRL has been widely adopted for HVAC control, building operation, distribution-network control, and renewable-building operation. These applications are primarily operational control problems in which states evolve through time. Static urban morphology optimization differs from that setting because each query is an absolute candidate description evaluated by a surrogate response surface.
\revisionseparator
\revisionloc{Main manuscript, Discussion paragraph.}
For early-stage urban energy design, the workflow is useful because it separates fast descriptor-space screening from feasible-block projection and sensitivity analysis. A surrogate optimizer can identify promising descriptor regions, while projection and selected physics-based calculations help determine how those regions relate to generated block geometries. This separation makes the workflow more transparent for design-support use.
\end{revisionbox}
\responseheading{Local references.}
\begin{responsereferences}
\referenceitem{Nguyen, A.-T., Reiter, S., \& Rigo, P. (2014). A review on simulation-based optimization methods applied to building performance analysis. \textit{Applied Energy}, 113, 1043--1058. \responsedoi{https://doi.org/10.1016/j.apenergy.2013.08.061}}
\referenceitem{Zhao, Z., Li, H., \& Wang, S. (2025). Surrogate-assisted coordinated design optimization of building and microclimate considering their mutual impacts. \textit{Applied Energy}, 383. \responsedoi{https://doi.org/10.1016/j.apenergy.2025.125374}}
\referenceitem{Ascione, F., De Masi, R. F., et al. (2016). Optimization of building envelope design for nZEBs in Mediterranean climate: Performance analysis of a residential case study. \textit{Applied Energy}, 183, 938--957. \responsedoi{https://doi.org/10.1016/j.apenergy.2016.09.027}}
\referenceitem{Du, Y., Zandi, H., et al. (2021). Intelligent multi-zone residential HVAC control strategy based on deep reinforcement learning. \textit{Applied Energy}, 281. \responsedoi{https://doi.org/10.1016/j.apenergy.2020.116117}}
\end{responsereferences}

\section*{Closing}
\addcontentsline{toc}{section}{Closing}
\responsepara{%
We again thank the editor and reviewers for the detailed comments. The revised manuscript and Supplementary Information now make the data-generation process, optimizer formulation, comparison protocol, physical and climate scope, and remaining limitations explicit. We hope that the revised package addresses the reviewers' concerns and provides a clearer basis for evaluating surrogate-conditioned optimizer comparisons in block-scale urban energy design.
}

\end{document}
